\documentclass{article}
\usepackage[utf8]{inputenc}
\usepackage{geometry}
\usepackage{enumitem}
\usepackage{hyperref}
\usepackage{xcolor}
\usepackage{graphicx}
\usepackage{array}
\usepackage{multirow}
\usepackage{amsmath, amssymb}
\usepackage{chemfig}
\usepackage{mhchem}
\usepackage{tikz}
\usepackage{setspace}
\usepackage{soul}
\usepackage{mathtools}
\usepackage{booktabs}
\usepackage{chemformula}
\usepackage{siunitx}
\usetikzlibrary{positioning, shapes, arrows.meta, decorations.pathreplacing, fit, backgrounds}
\usepackage{longtable}
\geometry{margin=1in}
\setlength{\parskip}{0.5em}
\usepackage{cancel}

\begin{document}

\begin{center}
    {\LARGE\textbf{SI Session Planning Form}}
\end{center}

\vspace{0.5cm}

\noindent
\begin{flushleft}
\textbf{SI Leader:} Brandon Cobb\\
\textbf{Session Week:} 3\\
\textbf{Session\#:} 3\\
\textbf{Instructor:} Professor Tramontozzi\\
\textbf{Old Plans:}\href{https://macombedu.sharepoint.com/:b:/s/ASC-SupplementalInstructionEmbeddedTutoring/EeSYqCwT_KhMolgJvXbVEjgBOyF_DevvS8BgCm1zPkQQJg?e=QJTt4o}{SharePoint}\\
\textbf{Session Goal:} Students will have reviewed unit conversions, algebraic manipulation of exact and inexact numbers, calculating their significant figures after multi-step operations like addition and multiplication. Student will review how to perform measurements using measurment devices with fixed uncertainty values, record data with units and write scienticific notation. Students are given a reference sheet for conversion factors.
\end{flushleft}

\vspace{1em}
\section*{\underline{Opener}}
{\centering\large\textbf{\hl{Take Attendance}}\par}\vspace{-0.25\baselineskip}
\noindent
\begin{tabular}{|p{4.2cm}|p{9.8cm}|}
\hline
\textbf{Collaborative Learning Techniques} & Clusters\\
\hline
\textbf{Learning Strategy} & Send a Problem\\
\hline
\textbf{Time} & 12:00---12:10\\
\hline
\textbf{Materials} & A writing utensil and index cards. \\
\hline
\textbf{Lecture/Textbook/Old Plans/Slide References} & Class roster, Chapter 1: Basics of Matter, Chapter 2: Measurments\\
\hline
\end{tabular}
\\
\vspace{0.5cm}
{\large\textbf{Definitions}}
\\
\begin{tabular}{|p{0.9\linewidth}|}
\hline
\begin{minipage}[t]{0.9\linewidth}
\begin{enumerate}
   \item\textbf{Clusters:} In Clusters, group participants are divided into smaller groups for discussion. They may also be allowed to self-select the small group they want to be in. After discussing the assigned topic, the cluster may report their findings to the large group.
   \item\textbf{Send a Problem:} Generate a list of problems. Assign a different problem to each student. Give the students a minute to complete the first step of their assigned problem. After the minute, students pass their problem to the students to their right, who then complete the second step of the problem. Continue passing the problem until the steps are completed. After, have the students present their originally-assigned problems with the completed solutions to the entire group.
\end{enumerate}
\end{minipage} //
\hline
\end{tabular}
\newpage
\noindent
{\large\textbf{Instructions (please provide a\hl{DETAILED} activity breakdown below (and/or attached}}\\

\begin{tabular}{|p{0.9\linewidth}|}
\hline
\vspace{0.2cm}
\textbf{Send a Problem (Chemistry Ice Breaker):}
Generate a list of chemistry problems of varying difficulty. Assign each student a different problem to solve. Give students one minute to complete the first step of their assigned problem. After the minute, each student passes their problem to the group on their right. The receiving group then completes the next step of the problem. This continues until all steps are completed. Finally, each student presents their originally-assigned problem along with the completed solution to the entire class.

\begin{enumerate}
\item The SI leader pairs up students based on proximity seating. If a student does not have a partner, the leader joins them.
\item Each pair is assigned a small chemistry task (e.g., determining concentration, calculating molar mass, predicting reaction products) suitable for a first step in a multi-step problem.
\item Pairs have two minutes to decide on their approach and collect any necessary tools or reference materials provided by the leader.
\item Pairs complete the first step of their problem, recording results clearly, and then pass the problem and any physical materials (if applicable) to the next group.
\item The next group evaluates the previous step and completes the subsequent step. If a group identifies an incorrect approach or calculation, they are allowed to correct it for the next step.
\item After all steps are completed, each original problem owner presents the full, completed solution to the class, highlighting any corrections or insights from peer contributions.
\item The SI leader facilitates discussion on problem-solving strategies, peer collaboration, and common chemistry misconceptions.
\end{enumerate}

Students engage collaboratively, seeing multiple approaches to the same problem, learning from their peers, and practicing stepwise problem-solving on a shared “centralized canvas” of the problem.\\
\hline
\end{tabular}

\newpage
\noindent
{\large\textbf{Content (fully worked out problems \& answers)}\\[0.2em]
\begin{tabular}{|p{0.9\linewidth}|}
\hline
\vspace{0.2cm}
\textbf{(1)} A chemist pours 0.375 L into a beaker already containing 125 mL. Total volume in liters? \\
\begin{minipage}[t]{0.85\linewidth}
\begin{enumerate}[label=\alph*.]
\item 0.500 L
\item 0.250 L
\item 0.375 L
\item 0.125 L
\end{enumerate}
\end{minipage}\\
\hline
\vspace{0.2cm}
\textbf{Solution:} \\

\[
125~\text{mL} \times \frac{1~\text{L}}{1000~\text{mL}} = 0.125~\text{L}
\] \\
\[
V_\text{total} = 0.375~\text{L} + 0.125~\text{L} = 0.500~\text{L}
\] \\ 
\hline
\textbf{Answer:} a. 0.500 L \\
\hline
\end{tabular}

\newpage
\noindent
{\large\textbf{Content (fully worked out problems \& answers) continued...}}\\[0.2em]
\vspace{0.2cm}
\begin{tabular}{|p{0.9\linewidth}|}
\hline
\vspace{0.2cm}
\textbf{(2)} Block of mass 540 g, volume half of 10. cm cube. Density? \\
\begin{minipage}[t]{0.85\linewidth}
\begin{enumerate}[label=\alph*.]
   \item 1.1 g/cm$^3$
   \item 0.54 g/cm$^3$
   \item 2.0 g/cm$^3$
   \item 1.08 g/cm$^3$
\end{enumerate}
\end{minipage}\\
\hline
\vspace{0.2cm}
\textbf{Solution:} \\
\[
V_\text{cube} = 10^3 = 1000~\text{cm}^3\] \\
\[
V_\text{metal} = \frac{1}{2} \times 1000 = 500~\text{cm}^3
\] \\
\[
\rho = \frac{m}{V} = \frac{540~\text{g}}{500~\text{cm}^3} = 1.1~\text{g/cm}^3
\] \\ 
\hline
\vspace{0.2cm}
\textbf{Answer:} a. 1.1 g/cm$^3$\\
\hline
\end{tabular}

\newpage
\noindent
{\large\textbf{Content (fully worked out problems \& answers) continued...}}\\[0.2em]
\vspace{0.2cm}
\begin{tabular}{|p{0.9\linewidth}|}
\hline
\vspace{0.2cm}
\textbf{(3)} 1.2 kg water poured into 2 equal beakers. Mass per beaker? \\
\begin{minipage}[t]{0.85\linewidth}
\begin{enumerate}[label=\alph*.]
   \item $6.0\times10^2$ g
   \item 6000 g
   \item 1.2 g
   \item $6.0\times10^1$ g
\end{enumerate}
\end{minipage}\\
\hline
\vspace{0.2cm}
\textbf{Solution:} \\
\[
1.2~\text{kg} \times 1000~\frac{\text{g}}{\text{kg}} = 1200~\text{g}
\] \\
\[
m_\text{per beaker} = \frac{1200~\text{g}}{2} = 6.0 \times 10^2~\text{g}
\] \\
\\
\hline
\vspace{0.2cm}
\textbf{Answer:} a. $6.0\times10^2$ g \\
\hline
\end{tabular}

\newpage
\noindent
{\large\textbf{Content (fully worked out problems \& answers) continued...}}\\[0.2em]
\vspace{0.2cm}
\begin{tabular}{|p{0.9\linewidth}|}
\hline
\vspace{0.2cm}
\textbf{(4)} Oil of 170 g, density 0.85 g/mL. Volume in liters? \\
\begin{minipage}[t]{0.85\linewidth}
\begin{enumerate}[label=\alph*.]
   \item 0.20 L
   \item 0.170 L
   \item 2.00 L
   \item 0.020 L
\end{enumerate}
\end{minipage}\\
\hline
\vspace{0.2cm}
\textbf{Solution:} \\
\[
V = \frac{m}{\rho} = \frac{170~\text{g}}{0.85~\text{g/mL}} = 200~\text{mL}
\] \\
\[
V_\text{L} = 200~\text{mL} \times \frac{1~\text{L}}{1000~\text{mL}} = 0.20~\text{L}
\] \\ 
\hline
\vspace{0.2cm}
\textbf{Answer:} a. 0.20 L \\
\hline
\end{tabular}

\newpage
\noindent
{\large\textbf{Content (fully worked out problems \& answers) continued...}}\\[0.2em]
\vspace{0.2cm}
\begin{tabular}{|p{0.9\linewidth}|}
\hline
\vspace{0.2cm}
\textbf{(5)} Cylinder 250 mL alcohol; 0.30 kg water added. Final volume in liters? \\
\begin{minipage}[t]{0.85\linewidth}
\begin{enumerate}[label=\alph*.]
   \item 0.55 L
   \item 0.35 L
   \item 0.75 L
   \item 0.25 L
\end{enumerate}
\end{minipage}\\
\hline
\vspace{0.2cm}
\textbf{Solution:} \\
\[
0.30~\text{kg} \times 1000~\frac{\text{g}}{\text{kg}} = 300~\text{g}
\] \\
\[
V_\text{water} = \frac{300~\text{g}}{1~\text{g/mL}} = 300~\text{mL}
\] \\
\[
V_\text{total} = 250~\text{mL} + 300~\text{mL} = 550~\text{mL}
\] \\
\[
V_\text{L} = 550~\text{mL} \times \frac{1~\text{L}}{1000~\text{mL}} = 0.55~\text{L}
\] \\ 
\hline
\vspace{0.2cm}
\textbf{Answer:} a. 0.55 L \\
\hline
\end{tabular}

\newpage
\noindent
{\large\textbf{Content (fully worked out problems \& answers) continued...}}\\[0.2em]
\vspace{0.2cm}
\begin{tabular}{|p{0.9\linewidth}|}
\hline
\vspace{0.2cm}
\textbf{(6)} Beaker contains 0.5 kg, half the required mass. Required mass in grams? \\
\begin{minipage}[t]{0.85\linewidth}
\begin{enumerate}[label=\alph*.]
   \item 1000 g
   \item 500 g
   \item 250 g
   \item 2000 g
\end{enumerate}
\end{minipage}\\
\hline
\vspace{0.2cm}
\textbf{Solution:} \\
\[
m_\text{required} = 0.5~\text{kg} \times 2 = 1.0~\text{kg}
\] \\
\[
1.0~\text{kg} \times 1000~\frac{\text{g}}{\text{kg}} = 1000~\text{g}
\] \\ 
\hline
\vspace{0.2cm}
\textbf{Answer:} a. 1000 g \\
\hline
\end{tabular}

\newpage
\noindent
\newpage
\noindent
{\large\textbf{Content (fully worked out problems& answers) continued...}}\\[0.2em]
\vspace{0.2cm}
\begin{tabular}{|p{0.9\linewidth}|}
\hline
\vspace{0.2cm}
\textbf{(7)} A student measures the mass of a sample as 12.5 g, but the true mass is 0.125 kg. What is the correct mass in grams? \\
\begin{minipage}[t]{0.85\linewidth}
\begin{enumerate}[label=\alph*.]
   \item 125 g
   \item 12.5 g
   \item 0.125 g
   \item 1.25 g
\end{enumerate}
\end{minipage}\\
\hline
\vspace{0.2cm}
\textbf{Solution:} \\
\[
0.125~\text{kg} \times 1000~\frac{\text{g}}{\text{kg}} = 125~\text{g}
\] \\ 
\hline
\vspace{0.2cm}
\textbf{Answer:} a. 125 g \\
\hline
\end{tabular}

\newpage
\noindent
{\large\textbf{Content (fully worked out problems \& answers) continued...}}\\[0.2em]
\vspace{0.2cm}
\begin{tabular}{|p{0.9\linewidth}|}
\hline
\vspace{0.2cm}
\textbf{(8)} Rectangular prism: volume twice mass in cm$^3$, mass 120 g. Density? \\
\begin{minipage}[t]{0.85\linewidth}
\begin{enumerate}[label=\alph*.]
   \item 0.50 g/cm$^3$
   \item 1.00 g/cm$^3$
   \item 0.25 g/cm$^3$
   \item 2.00 g/cm$^3$
\end{enumerate}
\end{minipage}\\
\hline
\vspace{0.2cm}
\textbf{Solution:} \\
\[
V = 2 \times 120~\text{cm}^3 = 240~\text{cm}^3 
\] \\
\[
\rho = \frac{m}{V} = \frac{120~\text{g}}{240~\text{cm}^3} = 0.50~\text{g/cm}^3
\] \\ 
\hline
\vspace{0.2cm}
\textbf{Answer:} a. 0.50 g/cm$^3$ \\
\hline
\end{tabular}

\newpage
\noindent
{\large\textbf{Content (fully worked out problems \& answers) continued...}}\\[0.2em]
\vspace{0.2cm}
\begin{tabular}{|p{0.9\linewidth}|}
\hline
\vspace{0.2cm}
\textbf{(9)} Mix 300. mL (0.90 g/mL) with 200. mL (1.2 g/mL). Average density? \\
\begin{minipage}[t]{0.85\linewidth}
\begin{enumerate}[label=\alph*.]
   \item 1.0 g/mL
   \item 0.98 g/mL
   \item 1.10 g/mL
   \item 0.90 g/mL
\end{enumerate}
\end{minipage}\\
\hline
\vspace{0.2cm}
\textbf{Solution:} \\
\[
m_A = 300~\text{mL} \times 0.9~\text{g/mL} = 270~\text{g}
\] \\
\[
m_B = 200~\text{mL} \times 1.2~\text{g/mL} = 240~\text{g}
\] \\
\[
V_\text{total} = 300.~\text{mL} + 200.~\text{mL} = 500.~\text{mL}
\] \\
\[
\rho_\text{avg} = \frac{m_\text{total}}{V_\text{total}} = \frac{270 + 240}{500.} = 1.0~\text{g/mL} 
\] \\ 
\hline
\vspace{0.2cm}
\textbf{Answer:} a. 1.0 g/mL \\
\hline
\end{tabular}

\newpage
\noindent
{\large\textbf{Content (fully worked out problems \& answers) continued...}}\\[0.2em]
\vspace{0.2cm}
\begin{tabular}{|p{0.9\linewidth}|}
\hline
\vspace{0.2cm}
\textbf{(10)} Flask contains 0.6 L; student removes 200 mL. Remaining volume in liters? \\
\begin{minipage}[t]{0.85\linewidth}
\begin{enumerate}[label=\alph*.]
   \item 0.4 L
   \item 0.2 L
   \item 0.8 L
   \item 0.6 L
\end{enumerate}
\end{minipage}\\
\hline
\vspace{0.2cm}
\textbf{Solution:} \\
\[
200~\text{mL} \times \frac{1~\text{L}}{1000~\text{mL}} = 0.2~\text{L}
\] \\
\[
V_\text{remaining} = 0.6~\text{L} - 0.200~\text{L} = 0.4~\text{L}
\] \\ 
\hline
\vspace{0.2cm}
\textbf{Answer:} a. 0.4 L \\
\hline
\end{tabular}

\newpage
\noindent
{\large\textbf{Content (fully worked out problems \& answers) continued...}}\\[0.2em]
\vspace{0.2cm}
\begin{tabular}{|p{0.9\linewidth}|}
\hline
\vspace{0.2cm}
\textbf{(11)} Metal mass 300 g, scale reads 0.3 kg. Correct grams? \\
\begin{minipage}[t]{0.85\linewidth}
\begin{enumerate}[label=\alph*.]
   \item Yes, 300 g
   \item No, 30 g
   \item No, 3 g
   \item No, 3000 g
\end{enumerate}
\end{minipage}\\
\hline
\vspace{0.2cm}
\textbf{Solution:} \\
\[
0.3~\text{kg} \times 1000~\frac{\text{g}}{\text{kg}} = 300~\text{g}
\] \\ 
\hline
\vspace{0.2cm}
\textbf{Answer:} a. Yes, 300 g \\
\hline
\end{tabular}

\newpage
\noindent
{\large\textbf{Content (fully worked out problems \& answers) continued...}}\\[0.2em]
\vspace{0.2cm}
\begin{tabular}{|p{0.9\linewidth}|}
\hline
\vspace{0.2cm}
\textbf{(12)} Rectangular prism: 10 cm × 5 cm, mass 250 g, density 2.5 g/cm$^3$. Height? \\
\begin{minipage}[t]{0.85\linewidth}
\begin{enumerate}[label=\alph*.]
   \item 2 cm
   \item 5 cm
   \item 1 cm
   \item 10 cm
\end{enumerate}
\end{minipage}\\
\hline
\vspace{0.2cm}
\textbf{Solution:} \\
\[
V = \frac{m}{\rho} = \frac{250~\text{g}}{2.5~\text{g/cm}^3} = 100~\text{cm}^3
\] \\
\[
h = \frac{V}{l \times w} = \frac{100~\text{cm}^3}{10~\text{cm} \times 5~\text{cm}} = 2~\text{cm}
\] \\ 
\hline
\vspace{0.2cm}
\textbf{Answer:} a. 2 cm \\
\hline
\end{tabular}

\newpage
\noindent
{\large\textbf{Content (fully worked out problems \& answers) continued...}}\\[0.2em]
\vspace{0.2cm}
\begin{tabular}{|p{0.9\linewidth}|}
\hline
\vspace{0.2cm}
\textbf{(13)} 50 mL liquid (0.8 g/mL); 25 mL evaporates. Mass left? \\
\begin{minipage}[t]{0.85\linewidth}
\begin{enumerate}[label=\alph*.]
   \item 20 g
   \item 40 g
   \item 10 g
   \item 25 g
\end{enumerate}
\end{minipage}\\
\hline
\vspace{0.2cm}
\textbf{Solution:} \\
\[
V_\text{left} = 50~\text{mL} - 25~\text{mL} = 25~\text{mL}
\] \\
\[
m = \rho \times V_\text{left} = 0.8~\text{g/mL} \times 25~\text{mL} = 20~\text{g}
\] \\ 
\hline
\vspace{0.2cm}
\textbf{Answer:} a. 20 g \\
\hline
\end{tabular}

\newpage
\noindent
{\large\textbf{Content (fully worked out problems \& answers) continued...}}\\[0.2em]
\vspace{0.2cm}
\begin{tabular}{|p{0.9\linewidth}|}
\hline
\vspace{0.2cm}
\textbf{(14)} 2.000 moles HCl, molar mass 36.46 g/mol. Mass should be? \\
\begin{minipage}[t]{0.85\linewidth}
\begin{enumerate}[label=\alph*.]
   \item 72.92 g
   \item 36.46 g
   \item 18.23 g
   \item 144.12 g
\end{enumerate}
\end{minipage}\\
\hline
\vspace{0.2cm}
\textbf{Solution:} \\
\[
m = n \times M = 2.000~\text{mol} \times 36.46~\text{g/mol} = 72.92~\text{g}
\] \\ 
\hline
\textbf{Answer:} a. 72.92 g \\
\hline
\end{tabular}

\newpage
\noindent
{\large\textbf{Content (fully worked out problems \& answers) continued...}}\\[0.2em]
\vspace{0.2cm}
\begin{tabular}{|p{0.9\linewidth}|}
\hline
\vspace{0.2cm}
\textbf{(15)} 0.5 L NaOH, add 250 mL, container reads 0.8 L. Consistent? \\
\begin{minipage}[t]{0.85\linewidth}
\begin{enumerate}[label=\alph*.]
   \item Yes
   \item No, it should read 0.7 L
   \item No, it should read 0.75 L
   \item No, it should read 1.0 L
\end{enumerate}
\end{minipage}\\
\hline
\vspace{0.2cm}
\textbf{Solution:} \\
\[
0.250~\text{L} + 0.5~\text{L} = 0.75~\text{L} 
\] \\ 
\hline
\vspace{0.2cm}
\textbf{Answer:} a. Yes \\
\hline
\end{tabular}

\newpage
\noindent
{\large\textbf{Content (fully worked out problems \& answers) continued...}}\\[0.2em]
\vspace{0.2cm}
\begin{tabular}{|p{0.9\linewidth}|}
\hline
\vspace{0.2cm}
\textbf{(16)} Cube of ice edge 4.00 cm, density 0.920 g/cm$^3$. Mass? \\
\begin{minipage}[t]{0.85\linewidth}
\begin{enumerate}[label=\alph*.]
   \item 58.9 g
   \item 64.0 g
   \item 56.0 g
   \item 60.0 g
\end{enumerate}
\end{minipage}\\
\hline
\vspace{0.2cm}
\textbf{Solution:} \\
\[
V = 4.00^3~\text{cm}^3 = 64.0~\text{cm}^3
\] \\
\[
m = \rho \times V = 0.920~\text{g/cm}^3 \times 64.0~\text{cm}^3 = 58.88~\text{g} \approx 58.9~\text{g}
\] \\ 
\hline
\vspace{0.2cm}
\textbf{Answer:} a. 58.9 g \\
\hline
\end{tabular}

\newpage
\noindent
{\large\textbf{Content (fully worked out problems \& answers) continued...}}\\[0.2em]
\vspace{0.2cm}
\begin{tabular}{|p{0.9\linewidth}|}
\hline
\vspace{0.2cm}
\textbf{(17)} Sample contains 0.03 mol Mg (24.31 g/mol). Mass? \\
\begin{minipage}[t]{0.85\linewidth}
\begin{enumerate}[label=\alph*.]
   \item 0.73 g
   \item 7.29 g
   \item 0.073 g
   \item 0.30 g
\end{enumerate}
\end{minipage}\\
\hline
\vspace{0.2cm}
\textbf{Solution:} \\
\[
m = n \times M = 0.03~\text{mol} \times 24.31~\text{g/mol} = 0.7293~\text{g} \approx 0.73~\text{g}
\] \\ 
\hline
\vspace{0.2cm}
\textbf{Answer:} a. 0.73 g \\
\hline
\end{tabular}

\newpage
\noindent
{\large\textbf{Content (fully worked out problems \& answers) continued...}}\\[0.2em]
\vspace{0.2cm}
\begin{tabular}{|p{0.9\linewidth}|}
\hline
\vspace{0.2cm}
\textbf{(18)} A student heats 30 g of ice. How many grams of water are formed? \\
\begin{minipage}[t]{0.85\linewidth}
\begin{enumerate}[label=\alph*.]
   \item 30 g
   \item 0 g
   \item 15 g
   \item 60 g
\end{enumerate}
\end{minipage}\\
\hline
\vspace{0.2cm}
\textbf{Solution:} \\
All the ice melts, so the mass of water is equal to the mass of ice:  
\[
\text{Mass of water} = 30~\text{g}
\] \\ 
\hline
\vspace{0.2cm}
\textbf{Answer:} a. 30 g \\
\hline
\end{tabular}


\newpage
\noindent
{\large\textbf{Content (fully worked out problems \& answers) continued...}}\\[0.2em]
\vspace{0.2cm}
\begin{tabular}{|p{0.9\linewidth}|}
\hline
\vspace{0.2cm}
\textbf{(19)} 0.900 kg ethanol, density 0.789 g/mL. Volume in mL?
\begin{minipage}[t]{0.85\linewidth}
\begin{enumerate}[label=\alph*.]
   \item 1140 mL
   \item 900 mL
   \item 700 mL
   \item 1200 mL
\end{enumerate}
\end{minipage}\\
\hline
\vspace{0.2cm}
\textbf{Solution:} \\

\[
0.900~\text{kg} \times 1000~\frac{\text{g}}{\text{kg}} = 900.~\text{g}
\] \\
\[
V = \frac{m}{\rho} = \frac{900.~\text{g}}{0.789~\text{g/mL}} \approx 1140~\text{mL}
\] \\ 
\hline
\vspace{0.2cm}
\textbf{Answer:} a. 1140 mL \\
\hline
\end{tabular}

\newpage
\noindent
{\large\textbf{Content (fully worked out problems \& answers) continued...}}\\[0.2em]
\vspace{0.2cm}
\begin{tabular}{|p{0.9\linewidth}|}
\hline
\vspace{0.2cm}
\textbf{(20)} Mix 150. mL (1.10 g/mL) with 50. mL (0.90 g/mL). Average density? \\
\begin{minipage}[t]{0.85\linewidth}
\begin{enumerate}[label=\alph*.]
   \item 1.05 g/mL
   \item 1.00 g/mL
   \item 0.95 g/mL
   \item 1.10 g/mL
\end{enumerate}
\end{minipage}\\
\hline
\vspace{0.2cm}
\textbf{Solution:} \\
\[
m_1 = 150.~\text{mL} \times 1.10~\text{g/mL} = 165~\text{g}, \quad
\] \\
\[
m_2 = 50.~\text{mL} \times 0.90~\text{g/mL} = 45~\text{g}
\] \\
\[
V_\text{total} = 150~\text{mL} + 50~\text{mL} = 200~\text{mL} 
\] \\
\[
\rho_\text{avg} = \frac{m_\text{total}}{V_\text{total}} = \frac{165 + 45}{200} = 1.05~\text{g/mL} 
\] \\ 
\hline
\vspace{0.2cm}
\textbf{Answer:} a. 1.05 g/mL \\
\hline
\end{tabular}

\newpage
\noindent
{\large\textbf{Content (fully worked out problems \& answers) continued...}}\\[0.2em]
\vspace{0.2cm}
\begin{tabular}{|p{0.9\linewidth}|}
\hline
\vspace{0.2cm}
\textbf{(21)} 300. mL water, 5\% evaporates. Remaining volume in mL and L? \\
\begin{minipage}[t]{0.85\linewidth}
\begin{enumerate}[label=\alph*.]
   \item 285 mL = 0.285 L
   \item 290 mL = 0.290 L
   \item 250 mL = 0.250 L
   \item 300 mL = 0.300 L
\end{enumerate}
\end{minipage}\\
\hline
\vspace{0.2cm}
\textbf{Solution:} \\
\[
V_\text{evap} = 0.05 \times 300.~\text{mL} = 15~\text{mL}
\] \\
\[
V_\text{left} = 300.~\text{mL} - 15~\text{mL} = 285~\text{mL}
\] \\
\[
V_\text{left} = 285~\text{mL} \times \frac{1~\text{L}}{1000~\text{mL}} = 0.285~\text{L}
\] \\ 
\hline
\vspace{0.2cm}
\textbf{Answer:} a. 285 mL = 0.285 L \\
\hline
\end{tabular}


\newpage
\noindent
{\large\textbf{Content (fully worked out problems \& answers) continued...}\\[0.2em]
\vspace{0.2cm}
\begin{tabular}{|p{0.9\linewidth}|}
\hline
\vspace{0.2cm}
\textbf{(22)} Block mass 200 g, density 0.8 g/cm$^3$. Student thinks volume 300 cm$^3$. Correct volume? \\
\begin{minipage}[t]{0.85\linewidth}
\begin{enumerate}[label=\alph*.]
   \item 300 cm$^3$
   \item 250 cm$^3$
   \item 200 cm$^3$
   \item 400 cm$^3$
\end{enumerate}
\end{minipage}\\
\hline
\vspace{0.2cm}
\textbf{Solution:} \\
\[
V = \frac{m}{\rho} = \frac{200~\text{g}}{0.8~\text{g/cm}^3} = 300~\text{cm}^3
\] \\ 
\hline
\vspace{0.2cm}
\textbf{Answer:} a. 300 cm$^3$ \\
\hline
\end{tabular}

\newpage
\noindent
{\large\textbf{Content (fully worked out problems& answers) continued...}\\[0.2em]
\vspace{0.2cm}
\begin{tabular}{|p{0.9\linewidth}|}
\hline
\vspace{0.2cm}
\textbf{(23)} A chemist adds 58.44 g of NaCl to a beaker containing 500. g of water. What is the total mass of the solution?
\begin{minipage}[t]{0.85\linewidth}
\begin{enumerate}[label=\alph*.]
   \item 558.44 g
   \item 542.00 g
   \item 500.00 g
   \item 58.44 g
\end{enumerate}
\end{minipage}\\
\hline
\vspace{0.2cm}
\textbf{Solution:} \\

\[
\text{Total mass} = \text{mass of solute} + \text{mass of solvent} = 58.44~\text{g} + 500.~\text{g} = 558.44~\text{g}
\] \\ 
\hline
\vspace{0.2cm}
\textbf{Answer:} a. 558.44 g \\
\hline
\end{tabular}  

\newpage
\noindent
{\large\textbf{Content (fully worked out problems& answers) continued...}\\[0.2em]
\vspace{0.2cm}
\begin{tabular}{|p{0.9\linewidth}|}
\hline
\vspace{0.2cm}
\textbf{(24)} A student measures a sugar sample as 120 mg instead of 0.12 g. What is the correct mass in grams? Explain \\
\begin{minipage}[t]{0.85\linewidth}
\begin{enumerate}[label=\alph*.]
   \item 0.12 g
   \item 1.20 g
   \item 0.012 g
   \item 12.0 g
\end{enumerate}
\end{minipage}\\
\hline
\vspace{0.2cm}
\textbf{Solution:} \\

\[
120~\text{mg} \times \frac{1~\text{g}}{1000~\text{mg}} = 0.12~\text{g} 
\] \\ 
\hline
\vspace{0.2cm}
\textbf{Answer:} a. 0.12 g, the answer must explain the conversion of 120 mg to 0.12 g, instead of assuming it is equal to 0.12 g. \\
\hline
\end{tabular}

\newpage
\noindent
{\large\textbf{Content (fully worked out problems \& answers) continued...}\\[0.2em]
\vspace{0.2cm}
% Q1
\begin{tabular}{|p{0.9\linewidth}|}
\hline
\vspace{0.2cm}
\textbf{(25)} What is my name?\\
\begin{minipage}[t]{0.85\linewidth}
\begin{enumerate}[label=\alph*.]
   \item Brandon
   \item Michael
   \item James
   \item Daniel
\end{enumerate}
\end{minipage}\\
\hline
\vspace{0.2cm}
\textbf{Answer:} a. Brandon\\
\hline
\end{tabular}
\newpage

% Q2

\noindent
{\large\textbf{Content (full worked out problems \& answers) continued...}\\[0.2em]
\vspace{0.2cm}
\begin{tabular}{|p{0.9\linewidth}|}
\hline
\vspace{0.2cm}
\textbf{(26)} What is the professor's name?\\
\begin{minipage}[t]{0.85\linewidth}
\begin{enumerate}[label=\alph*.]
   \item Dr. Smith
   \item Dr. Johnson
   \item Dr. T, Dr. Tramontozzi
   \item Dr. Brown
\end{enumerate}
\end{minipage}\\
\hline
\vspace{0.2cm}
\textbf{Answer:} c. Dr. T, Dr. Tramontozzi\\
\hline
\end{tabular}
\newpage

% Q3

\noindent
{\large\textbf{Content (full worked out problems \& answers) continued...}\\[0.2em]
\vspace{0.2cm}
\begin{tabular}{|p{0.9\linewidth}|}
\hline
\vspace{0.2cm}
\textbf{(27)} What days is class and lab held?\\
\begin{minipage}[t]{0.85\linewidth}
\begin{enumerate}[label=\alph*.]
   \item Monday and Wednesday
   \item Tuesday and Thursday
   \item Friday only
   \item Wednesday and Friday
\end{enumerate}
\end{minipage}\\
\hline
\vspace{0.2cm}
\textbf{Answer:} b. Tuesday and Thursday\\
\hline
\end{tabular}
\newpage

% Q4

\noindent
{\large\textbf{Content (full worked out problems \& answers) continued...}\\[0.2em]
\vspace{0.2cm}
\begin{tabular}{|p{0.9\linewidth}|}
\hline
\vspace{0.2cm}
\textbf{(28)} What time does class start and when does lab end?\\
\begin{minipage}[t]{0.85\linewidth}
\begin{enumerate}[label=\alph*.]
   \item 9:00 AM – 12:00 PM
   \item 8:00 AM – 11:30 AM
   \item 10:00 AM – 1:30 PM
   \item 1:00 PM – 4:30 PM
\end{enumerate}
\end{minipage}\\
\hline
\vspace{0.2cm}
\textbf{Answer:} b. 8:00 AM – 11:30 AM\\
\hline
\end{tabular}
\newpage

% Q5

\noindent
{\large\textbf{Content (full worked out problems \& answers) continued...}\\[0.2em]
\vspace{0.2cm}
\begin{tabular}{|p{0.9\linewidth}|}
\hline
\vspace{0.2cm}
\textbf{(29)} What days are the supplemental instruction sessions?\\
\begin{minipage}[t]{0.85\linewidth}
\begin{enumerate}[label=\alph*.]
   \item Mondays
   \item Tuesdays
   \item Fridays
   \item Thursdays
\end{enumerate}
\end{minipage}\\
\hline
\vspace{0.2cm}
\textbf{Answer:} b. Tuesdays\\
\hline
\end{tabular}
\newpage

% Q6

\noindent
{\large\textbf{Content (full worked out problems \& answers) continued...}\\[0.2em]
\vspace{0.2cm}
\begin{tabular}{|p{0.9\linewidth}|}
\hline
\vspace{0.2cm}
\textbf{(30)} What times are the supplemental instruction sessions?\\
\begin{minipage}[t]{0.85\linewidth}
\begin{enumerate}[label=\alph*.]
   \item 9:00 AM – 10:00 AM
   \item 11:50 AM – 12:50 PM
   \item 2:00 PM – 3:00 PM
   \item 4:00 PM – 5:00 PM
\end{enumerate}
\end{minipage}\\
\hline
\vspace{0.2cm}
\textbf{Answer:} b. 11:50 AM – 12:50 PM\\
\hline
\end{tabular}
\newpage

% Q7

\noindent
{\large\textbf{Content (full worked out problems \& answers) continued...}\\[0.2em]
\vspace{0.2cm}
\begin{tabular}{|p{0.9\linewidth}|}
\hline
\textbf{(31)} What is OWL?\\
\vspace{0.2cm}
\begin{minipage}[t]{0.85\linewidth}
\begin{enumerate}[label=\alph*.]
   \item A campus tutoring service
   \item Online Web Learning platform used for assignments and course resources
   \item A student club
   \item A type of exam
\end{enumerate}
\end{minipage}\\
\hline
\vspace{0.2cm}
\textbf{Answer:} b. Online Web Learning platform used for assignments and course resources\\
\hline
\end{tabular}
\newpage

% Q8

\noindent
{\large\textbf{Content (full worked out problems \& answers) continued...}\\[0.2em]
\vspace{0.2cm}
\begin{tabular}{|p{0.9\linewidth}|}
\hline
\vspace{0.2cm}
\textbf{(32)} What is Canvas?\\
\begin{minipage}[t]{0.85\linewidth}
\begin{enumerate}[label=\alph*.]
   \item A classroom whiteboard
   \item A learning management system where grades and course materials are posted
   \item A chemistry textbook
   \item A scheduling app
\end{enumerate}
\end{minipage}\\
\hline
\vspace{0.2cm}
\textbf{Answer:} b. A learning management system where grades and course materials are posted\\
\hline
\end{tabular}
\newpage

% Q9

\noindent
{\large\textbf{Content (full worked out problems \& answers) continued...}\\[0.2em]
\vspace{0.2cm}
\begin{tabular}{|p{0.9\linewidth}|}
\hline
\vspace{0.2cm}
\textbf{(33)} What are the professor's office hours?\\
\begin{minipage}[t]{0.85\linewidth}
\begin{enumerate}[label=\alph*.]
   \item Monday 10–11 AM only
   \item Tuesday 11:30 AM – 12 PM; Wednesday 7:30–8 AM, 11:30–2 PM, 5–6 PM; Thursday 11:30 AM – 12 PM
   \item Friday 1–3 PM only
   \item By appointment only
\end{enumerate}
\end{minipage}\\
\hline
\vspace{0.2cm}
\textbf{Answer:} b. Tuesday 11:30 AM – 12 PM; Wednesday 7:30–8 AM, 11:30–2 PM, 5–6 PM; Thursday 11:30 AM – 12 PM\\
\hline
\end{tabular}
\newpage

\noindent
{\large\textbf{Content (full worked out problems \& answers) continued...}\\[0.2em]
\vspace{0.2cm}
\begin{tabular}{|p{0.9\linewidth}|}
\hline
\vspace{0.2cm}
\textbf{(34)} Which room is the class in?\\
\begin{minipage}[t]{0.85\linewidth}
\begin{enumerate}[label=\alph*.]
   \item SB 210
   \item SB 208
   \item SB 105
   \item SB 416
\end{enumerate}
\end{minipage}\\
\hline
\vspace{0.2cm}
\textbf{Answer:} a. SB 210\\
\hline
\end{tabular}

\newpage
\noindent
{\large\textbf{Content (full worked out problems \& answers) continued...}\\[0.2em]
\vspace{0.2cm}
\begin{tabular}{|p{0.9\linewidth}|}
\hline
\vspace{0.2cm}
\textbf{(35)} Which room is the lab in?\\
\begin{minipage}[t]{0.85\linewidth}
\begin{enumerate}[label=\alph*.]
   \item SB 416
   \item SB 210
   \item SB 105
   \item SB 310
\end{enumerate}
\end{minipage}\\
\hline
\vspace{0.2cm}
\textbf{Answer:} a. SB 416\\
\hline
\end{tabular}

\newpage
\noindent
{\large\textbf{Content (full worked out problems \& answers) continued...}\\[0.2em]
\begin{tabular}{|p{0.9\linewidth}|}
\hline
\vspace{0.2cm}
\textbf{(36)} Which room is the supplemental instruction in?\\
\begin{minipage}[t]{0.85\linewidth}
\begin{enumerate}[label=\alph*.]
   \item SB 208
   \item SB 210
   \item SB 105
   \item SB 416
\end{enumerate}
\end{minipage}\\
\hline
\vspace{0.2cm}
\textbf{Answer:} a. SB 208\\
\hline
\end{tabular}

\newpage
\noindent
{\large\textbf{Content (full worked out problems \& answers) continued...}\\[0.2em]
\begin{tabular}{|p{0.9\linewidth}|}
\hline
\vspace{0.2cm}
\textbf{(37)} Which campus are we at?\\
\begin{minipage}[t]{0.85\linewidth}
\begin{enumerate}[label=\alph*.]
   \item North Campus
   \item South Campus
   \item East Campus
   \item West Campus
\end{enumerate}
\end{minipage}\\
\hline
\vspace{0.2cm}
\textbf{Answer:} b. South Campus\\
\hline
\end{tabular}

\newpage
\noindent
{\large\textbf{Content (full worked out problems \& answers) continued...}\\[0.2em]
\begin{tabular}{|p{0.9\linewidth}|}
\hline
\textbf{(38)} What does a ruler measure?\\
\vspace{0.2cm}
\begin{minipage}[t]{0.85\linewidth}
\begin{enumerate}[label=\alph*.]
   \item Mass
   \item Length
   \item Temperature
   \item Volume
\end{enumerate}
\end{minipage}\\
\hline
\vspace{0.2cm}
\textbf{Answer:} b. Length\\
\hline
\end{tabular}

\newpage
\noindent
{\large\textbf{Content (full worked out problems \& answers) continued...}\\[0.2em]
\vspace{0.2cm}
\begin{tabular}{|p{0.9\linewidth}|}
\hline
\vspace{0.2cm}
\textbf{(39)} What does a scale measure?\\
\begin{minipage}[t]{0.85\linewidth}
\begin{enumerate}[label=\alph*.]
   \item Mass
   \item Temperature
   \item Length
   \item Pressure
\end{enumerate}
\end{minipage}\\
\hline
\vspace{0.2cm}
\textbf{Answer:} a. Mass\\
\hline
\end{tabular}
\newpage
\noindent
{\large\textbf{Content (full worked out problems \& answers) continued...}\\[0.2em]
\vspace{0.2cm}
\begin{tabular}{|p{0.9\linewidth}|}
\hline
\vspace{0.2cm}
\textbf{(40)} What does a volumetric flask measure?\\
\begin{minipage}[t]{0.85\linewidth}
\begin{enumerate}[label=\alph*.]
   \item Approximate volume
   \item Temperature
   \item Precise volume of liquids
   \item Mass
\end{enumerate}
\end{minipage}\\
\hline
\vspace{0.2cm}
\textbf{Answer:} c. Precise volume of liquids\\
\hline
\end{tabular}

\newpage
\noindent
{\large\textbf{Content (full worked out problems \& answers) continued...}\\[0.2em]
\vspace{0.2cm}
\begin{tabular}{|p{0.9\linewidth}|}
\hline
\vspace{0.2cm}
\textbf{(41)} What does a thermometer measure?\\
\begin{minipage}[t]{0.85\linewidth}
\begin{enumerate}[label=\alph*.]
   \item Pressure
   \item Temperature
   \item Volume
   \item Mass
\end{enumerate}
\end{minipage}\\
\hline
\vspace{0.2cm}
\textbf{Answer:} b. Temperature\\
\hline
\end{tabular}

\newpage
\noindent
{\large\textbf{Content (full worked out problems \& answers) continued...}\\[0.2em]
\vspace{0.2cm}
\begin{tabular}{|p{0.9\linewidth}|}
\hline
\vspace{0.2cm}
\textbf{(42)} What does a graduated cylinder measure?\\
\begin{minipage}[t]{0.85\linewidth}
\begin{enumerate}[label=\alph*.]
   \item Mass
   \item Volume of liquids
   \item Temperature
   \item Pressure
\end{enumerate}
\end{minipage}\\
\hline
\vspace{0.2cm}
\textbf{Answer:} b. Volume of liquids\\
\hline
\end{tabular}

\newpage
\noindent
{\large\textbf{Content (full worked out problems \& answers) continued...}\\[0.2em]
\vspace{0.2cm}
\begin{tabular}{|p{0.9\linewidth}|}
\hline
\vspace{0.2cm}
\textbf{(43)} What does a pipette measure?\\
\begin{minipage}[t]{0.85\linewidth}
\begin{enumerate}[label=\alph*.]
   \item Approximate volume
   \item Small, precise amounts of liquid
   \item Mass
   \item Temperature
\end{enumerate}
\end{minipage}\\
\hline
\vspace{0.2cm}
\textbf{Answer:} b. Small, precise amounts of liquid\\
\hline
\end{tabular}

\newpage
\noindent
{\large\textbf{Content (full worked out problems \& answers) continued...}\\[0.2em]
\vspace{0.2cm}
\begin{tabular}{|p{0.9\linewidth}|}
\hline
\vspace{0.2cm}
\textbf{(44)} What does a beaker measure?\\
\begin{minipage}[t]{0.85\linewidth}
\begin{enumerate}[label=\alph*.]
   \item Exact volume of liquids
   \item Approximate volume of liquids
   \item Temperature
   \item Mass
\end{enumerate}
\end{minipage}\\
\hline
\vspace{0.2cm}
\textbf{Answer:} b. Approximate volume of liquids\\
\hline
\end{tabular}
\newpage
\noindent
{\large\textbf{Content (full worked out problems \& answers) continued...}\\[0.2em]
\vspace{0.2cm}
\begin{tabular}{|p{0.9\linewidth}|}
\hline
\vspace{0.2cm}
\textbf{(45)} What does a balance measure?\\
\begin{minipage}[t]{0.85\linewidth}
\begin{enumerate}[label=\alph*.]
   \item Volume
   \item Mass of objects or substances
   \item Temperature
   \item Pressure
\end{enumerate}
\end{minipage}\\
\hline
\vspace{0.2cm}
\textbf{Answer:} b. Mass of objects or substances\\
\hline
\end{tabular}

\newpage
\noindent
{\large\textbf{Content (full worked out problems \& answers) continued...}\\[0.2em]
\vspace{0.2cm}
\begin{tabular}{|p{0.9\linewidth}|}
\hline
\vspace{0.2cm}
\textbf{(46)} What are significant figures?\\
\begin{minipage}[t]{0.85\linewidth}
\begin{enumerate}[label=\alph*.]
   \item The digits in a number which follow strict rules about precision
   \item All numbers written in scientific notation
   \item Only the digits before the decimal point
   \item Only the zeroes in a measurement
\end{enumerate}
\end{minipage}\\
\hline
\vspace{0.2cm}
\textbf{Answer:} a. The digits in a number which follow strict rules about precision\\
\hline
\end{tabular}

\newpage
\noindent
{\large\textbf{Content (full worked out problems \& answers) continued...}\\[0.2em]
\vspace{0.2cm}
\begin{tabular}{|p{0.9\linewidth}|}
\hline
\vspace{0.2cm}
\textbf{(47)} What's the difference of significant figures betweenexact and inexact numbers?\\
\begin{minipage}[t]{0.85\linewidth}
\begin{enumerate}[label=\alph*.]
   \item Exact numbers have infinite significant figures; inexact depend on measurement device
   \item Exact numbers have no significant figures; inexact always have three
   \item Exact numbers are estimates; inexact are precise
   \item Exact numbers are rounded; inexact numbers are not
\end{enumerate}
\end{minipage}\\
\hline
\vspace{0.2cm}
\textbf{Answer:} a. Exact numbers have infinite significant figures; inexact depend on measurement device\\
\hline
\end{tabular}

\newpage
\noindent
{\large\textbf{Content (full worked out problems \& answers) continued...}\\[0.2em]
\vspace{0.2cm}
\begin{tabular}{|p{0.9\linewidth}|}
\hline
\vspace{0.2cm}
\textbf{(48)} What is the fastest achievable speed?\\
\begin{minipage}[t]{0.85\linewidth}
\begin{enumerate}[label=\alph*.]
   \item Speed of sound
   \item Speed of light
   \item Speed of electrons in an atom
   \item Escape velocity of Earth
\end{enumerate}
\end{minipage}\\
\hline
\vspace{0.2cm}
\textbf{Answer:} b. Speed of light, $c\approx 3.00\times 10^{8}\,\text{m/s}$\\
\hline
\end{tabular}

\newpage
\noindent
{\large\textbf{Content (full worked out problems \& answers) continued...}\\[0.2em]
\vspace{0.2cm}
\begin{tabular}{|p{0.9\linewidth}|}
\hline
\vspace{0.2cm}
\textbf{(49)} Which is more important in representing a number which is subtracted?\\
\begin{minipage}[t]{0.85\linewidth}
\begin{enumerate}[label=\alph*.]
   \item The answer must contain the same number of decimals of the lesser precise measurements.
   \item The answer must contain the same number of decimals of the more precise measurements.
   \item The answer must contain the same number of significant figures of the lesser precise measurement.
   \item The answer must contain the same number of significant figures of the more precise measurements.
\end{enumerate}
\end{minipage}\\
\hline
\vspace{0.2cm}
\textbf{Answer:} a. The answer must contain the same number of decimals of the lesser precise measurements.\\
\hline
\end{tabular}
\newpage 
\newpage
\noindent
{\large\textbf{Content (fully worked out problems \& answers) continued...}}\\[0.2em]
\begin{tabular}{|p{0.9\linewidth}|}
\hline
\vspace{0.2cm}
\textbf{(50)} Convert blood glucose level of 125.0 mg/dL to g/L.\\
\begin{minipage}[t]{0.85\linewidth}
\begin{enumerate}[label=\alph*.]
   \item 0.0125 g/L
   \item 0.125 g/L
   \item 1.25 g/L
   \item 12.5 g/L
\end{enumerate}
\end{minipage}\\
\hline
\vspace{0.2cm}
\textbf{Answer:}\\
\[
\text{Concentration} =\frac{125~\cancel{\text{mg}}}{1~\cancel{\text{dL}}}\times\frac{1~\text{g}}{1000~\cancel{\text{mg}}}\times\frac{10~\cancel{\text{dL}}}{1~\text{L}} = 0.125~\text{g/L}
\] \\
\hline
\end{tabular}

\newpage
\noindent
{\large\textbf{Content (fully worked out problems \& answers) continued...}}\\[0.2em]
\begin{tabular}{|p{0.9\linewidth}|}
\hline
\vspace{0.2cm}
\textbf{(51)} Convert 0.085 $\text{m}^3$ to gallons. (Use $1~\text{m}^3=35.3147~\text{ft}^3$, $1~\text{ft}^3=7.48052~\text{gal}$)\\
\begin{minipage}[t]{0.85\linewidth}
\begin{enumerate}[label=\alph*.]
   \item 2.25 gal
   \item 22.5 gal
   \item 225 gal
   \item 2250 gal
\end{enumerate}
\end{minipage}\\
\hline
\vspace{0.2cm}
\textbf{Answer:}\\
\[
\text{Volume} =\frac{0.0850~\cancel{\text{m}^3}}{1}\times\frac{35.3147~\text{ft}^3}{1~\cancel{\text{m}^3}}\times\frac{7.48052~\text{gal}}{1~\cancel{\text{ft}^3}} = 22.5~\text{gal}
\] \\
\hline
\end{tabular}

\newpage
\noindent
{\large\textbf{Content (fully worked out problems \& answers) continued...}}\\[0.2em]
\begin{tabular}{|p{0.9\linewidth}|}
\hline
\vspace{0.2cm}
\textbf{(52)} Lead has density $11.34~\text{g/cm}^3$. What is the mass in kg of a 275.0 $\text{cm}^3$ sample?\\
\begin{minipage}[t]{0.85\linewidth}
\begin{enumerate}[label=\alph*.]
   \item 0.312 kg
   \item 3.12 kg
   \item 31.2 kg
   \item 312 kg
\end{enumerate}
\end{minipage}\\
\hline
\vspace{0.2cm}
\textbf{Answer:}\\
\[
\frac{11.34~\text{g}}{1~\cancel{\text{cm}^3}}\times\frac{275~\cancel{\text{cm}^3}}{1} = 3.12\times 10^3~\text{g}
\]
\[
\text{Mass} =\frac{3.12\times 10^3~\cancel{\text{g}}}{1}\times\frac{1~\text{kg}}{1000~\cancel{\text{g}}} = 3.12~\text{kg}
\] \\
\hline
\end{tabular}
\newpage
\noindent
{\large\textbf{Content (fully worked out problems \& answers) continued...}}\\[0.2em]
\vspace{0.2cm}
\begin{tabular}{|p{0.9\linewidth}|}
\hline
\vspace{0.2cm}
\textbf{(53)} A tectonic plate moves 860 km per day. What is its speed in m/s?\\
\begin{minipage}[t]{0.85\linewidth}
\begin{enumerate}[label=\alph*.]
   \item 0.995 m/s
   \item 9.95 m/s
   \item 95.9 m/s
   \item 0.0959 m/s
\end{enumerate}
\end{minipage}\\
\hline
\vspace{0.2cm}
\textbf{Answer:}\\
\[
\text{Speed}=\frac{860~\cancel{\text{km}}}{\cancel{\text{day}}}\times\frac{1000~\text{m}}{1~\cancel{\text{km}}}\times\frac{1~\cancel{\text{day}}}{86400~\text{s}}=9.95~\text{m/s}
\] \\
\hline
\end{tabular}

\newpage
\noindent
{\large\textbf{Content (fully worked out problems \& answers) continued...}}\\[0.2em]
\vspace{0.2cm}
\begin{tabular}{|p{0.9\linewidth}|}
\hline
\vspace{0.2cm}
\textbf{(54)} Convert 750 mL to $\text{in}^3$. (Use $1~\text{in}=2.54~\text{cm}$)\\
\begin{minipage}[t]{0.85\linewidth}
\begin{enumerate}[label=\alph*.]
   \item 4.58 $\text{in}^3$
   \item 45.8 $\text{in}^3$
   \item 458 $\text{in}^3$
   \item 0.458 $\text{in}^3$
\end{enumerate}
\end{minipage}\\
\hline
\vspace{0.2cm}
\textbf{Answer:}\\
\[
750~\text{mL}=750~\text{cm}^3,\quad 1~\text{in}^3=(2.54~\text{cm})^3=16.387~\text{cm}^3
\]
\[
\text{Volume}=\frac{750~\text{cm}^3}{16.387~\text{cm}^3/\text{in}^3}=45.8~\text{in}^3
\] \\
\hline
\end{tabular}

\newpage
\noindent
{\large\textbf{Content (fully worked out problems \& answers) continued...}}\\[0.2em]
\vspace{0.2cm}
\begin{tabular}{|p{0.9\linewidth}|}
\hline
\vspace{0.2cm}
\textbf{(55)} Convert 88.0 kPa to atm. (Use $1~\text{atm}=101.325~\text{kPa}$)\\
\begin{minipage}[t]{0.85\linewidth}
\begin{enumerate}[label=\alph*.]
   \item 0.868 atm
   \item 0.880 atm
   \item 1.15 atm
   \item 0.815 atm
\end{enumerate}
\end{minipage}\\
\hline
\vspace{0.2cm}
\textbf{Answer:}\\
\[
\text{Pressure}=\frac{88.0~\cancel{\text{kPa}}}{1}\times\frac{1~\text{atm}}{101.325~\cancel{\text{kPa}}}=0.868~\text{atm}
\] \\
\hline
\end{tabular}

\newpage
\noindent
{\large\textbf{Content (fully worked out problems \& answers) continued...}}\\[0.2em]
\vspace{0.2cm}
\begin{tabular}{|p{0.9\linewidth}|}
\hline
\vspace{0.2cm}
\textbf{(56)} Convert 4.6 g/L to mg/dL.\\
\begin{minipage}[t]{0.85\linewidth}
\begin{enumerate}[label=\alph*.]
   \item 0.46 mg/dL
   \item 4.6 mg/dL
   \item 46 mg/dL
   \item 460 mg/dL
\end{enumerate}
\end{minipage}\\
\hline
\vspace{0.2cm}
\textbf{Answer:}\\
\[
\frac{4.6~\cancel{\text{g}}}{\cancel{\text{L}}}\times\frac{1000~\text{mg}}{1~\cancel{\text{g}}}\times\frac{1~\text{dL}}{10~\cancel{\text{L}}}=46~\text{mg/dL}
\] \\
\hline
\end{tabular}

\newpage
\noindent
{\large\textbf{Content (fully worked out problems \& answers) continued...}}\\[0.2em]
\vspace{0.2cm}
\begin{tabular}{|p{0.9\linewidth}|}
\hline
\vspace{0.2cm}
\textbf{(57)} A block has mass 245.0 g and dimensions 4.50 cm $\times$ 3.20 cm $\times$ 2.10 cm. Find its density in kg/m$^3$.\\
\begin{minipage}[t]{0.85\linewidth}
\begin{enumerate}[label=\alph*.]
   \item $8.10\times10^{2}~\text{kg/m}^3$
   \item $8.10\times10^{3}~\text{kg/m}^3$
   \item $8.10\times10^{4}~\text{kg/m}^3$
   \item $8.10\times10^{5}~\text{kg/m}^3$
\end{enumerate}
\end{minipage}\\
\hline
\vspace{0.2cm}
\textbf{Answer:}\\
\[
V=(4.50)(3.20)(2.10)~\text{cm}^3=30.2~\text{cm}^3
\]
\[
\rho=\frac{245.0~\text{g}}{30.2~\text{cm}^3}=8.10~\text{g/mL}=8.10~\frac{\text{g}}{\text{cm}^3}
\]
\[
\text{Density}=\left(8.10~\frac{\text{g}}{\text{cm}^3}\right)\times\frac{1~\text{kg}}{1000~\text{g}}\times\frac{10^6~\text{cm}^3}{1~\text{m}^3}=8.10\times10^{3}~\text{kg/m}^3
\] \\
\hline
\end{tabular}

\newpage
\noindent
{\large\textbf{Content (fully worked out problems \& answers) continued...}}\\[0.2em]
\vspace{0.2cm}
\begin{tabular}{|p{0.9\linewidth}|}
\hline
\vspace{0.2cm}
\textbf{(58)} Show that 125.0 N$\cdot$m equals 125.0 J using unit analysis. (Use $1~\text{N}=1~\text{kg}\cdot\text{m/s}^2$)\\
\begin{minipage}[t]{0.85\linewidth}
\begin{enumerate}[label=\alph*.]
   \item 125.0 J
   \item 125.0 N
   \item 125.0 W
   \item 125.0 kg$\cdot$m/s$^2$
\end{enumerate}
\end{minipage}\\
\hline
\vspace{0.2cm}
\textbf{Answer:}\\
\[
125.0~\text{N}\cdot\text{m}=125.0\left(\frac{\text{kg}\cdot\text{m}}{\text{s}^2}\right)\cdot\text{m}=125.0~\frac{\text{kg}\cdot\text{m}^2}{\text{s}^2}=125.0~\text{J}
\] \\
\hline
\end{tabular}

\newpage
\noindent
{\large\textbf{Content (fully worked out problems& answers) continued...}}\\
\vspace{0.2cm}
\begin{tabular}{|p{0.9\linewidth}|}
\hline
\vspace{0.2cm}
\textbf{(59)} A student weighs a sample of water and records 248.5 g. The actual mass is 250.0 g. What is the difference between the recorded mass and the true mass?\\
\hline
\vspace{0.2cm}
\textbf{Solution:} \\

\[
\text{Difference} = \text{True mass} - \text{Recorded mass} = 250.0~\text{g} - 248.5~\text{g} = 1.5~\text{g}
\] \\ 
\hline
\vspace{0.2cm}
\textbf{Answer:} 1.5 g \\
\hline
\end{tabular}

\newpage
\noindent
{\large\textbf{Content (fully worked out problems \& answers) continued...}}\\
\vspace{0.2cm}
\begin{tabular}{|p{0.9\linewidth}|}
\hline
\vspace{0.2cm}
\textbf{(60)} A 0.500 g sample of NaCl dissolves in 250.0 mL of water. What is the concentration in g/L?\\
\hline
\vspace{0.2cm}
\begin{minipage}[t]{0.85\linewidth}
\begin{enumerate}[label=\alph*.]
\item 0.125 g/L
\item 2.00 g/L
\item 1.25 g/L
\item 0.250 g/L
\end{enumerate}
\end{minipage} \\
\hline
\vspace{0.2cm}
\textbf{Answer:} 
\[
\text{Concentration} = \frac{0.500~\text{g}}{0.250~\text{L}} = 2.00~\text{g/L}
\] \\
\\
\hline
\end{tabular}

\newpage
\noindent
{\large\textbf{Content (fully worked out problems \& answers) continued...}}\\[0.2em]
\vspace{0.2cm}
\begin{tabular}{|p{0.9\linewidth}|}
\hline
\vspace{0.2cm}
\textbf{(61)} Convert 987.0 mbar to Pa.\\
\hline
\vspace{0.2cm}
\begin{minipage}[t]{0.85\linewidth}
\begin{enumerate}[label=\alph*.]
\item 9.870 Pa
\item 9.87 $\times$ 10$^2$ Pa
\item 9.87 $\times$ 10$^4$ Pa
\item 9.87 $\times$ 10$^5$ Pa
\end{enumerate}
\end{minipage} \\
\hline
\vspace{0.2cm}
\textbf{Answer:} 
\[
987.0~\text{mbar} \times \frac{100~\text{Pa}}{1~\text{mbar}} = 9.870 \times 10^4~\text{Pa}
\] \\ 
\hline
\end{tabular}


\newpage
\noindent
{\large\textbf{Content (fully worked out problems \& answers) continued...}}\\[0.2em]
\vspace{0.2cm}
\begin{tabular}{|p{0.9\linewidth}|}
\hline
\vspace{0.2cm}
\textbf{(62)} Water has density 62.4 lb/ft$^3$. Convert to g/mL. (Use \(1~\text{lb}=453.592~\text{g}, 1~\text{ft}^3=28316.8467~\text{mL}\))\\
\hline
\vspace{0.2cm}
\begin{minipage}[t]{0.85\linewidth}
\begin{enumerate}[label=\alph*.]
\item 0.624 g/mL
\item 1.00 g/mL
\item 2.20 g/mL
\item 6.24 g/mL
\end{enumerate}
\end{minipage}\\
\hline
\vspace{0.2cm}
\textbf{Answer:} 
\[
62.4~\text{lb/ft}^3 \times \frac{453.592~\text{g}}{1~\text{lb}} \times \frac{1~\text{ft}^3}{28316.8467~\text{mL}} = 1.00~\text{g/mL}
\] \\
\hline
\end{tabular}

\newpage
\noindent
{\large\textbf{Content (fully worked out problems \& answers) continued...}}\\[0.2em]
\vspace{0.2cm}
\begin{tabular}{|p{0.9\linewidth}|}
\hline
\vspace{0.2cm}
\textbf{(63)} A sample occupies 250.0 mL and has a mass of 0.890 g. Calculate its density.\\
\hline
\vspace{0.2cm}
\begin{minipage}[t]{0.85\linewidth}
\begin{enumerate}[label=\alph*.]
\item 3.56 \times $10^{-3}$ g/mL
\item 0.890 g/mL
\item 3.56 g/mL
\item 35.6 g/mL
\end{enumerate}
\end{minipage}\\
\hline
\vspace{0.2cm}
\textbf{Answer:} 
\[
\rho = \frac{0.890~\text{g}}{250.0~\text{mL}} = 3.56 \times 10^{-3}~\text{g/mL}
\] \\ 
\hline
\end{tabular}

\newpage
\noindent
{\large\textbf{Content (fully worked out problems \& answers) continued...}}\\[0.2em]
\vspace{0.2cm}
\begin{tabular}{|p{0.9\linewidth}|}
\hline
\vspace{0.2cm}
\textbf{(64)} A beaker contains 3.50 L of solution. Express this volume in cubic decimeters (dm$^3$) using unit factors.\\
\hline
\vspace{0.2cm}
\begin{minipage}[t]{0.85\linewidth}
\begin{enumerate}[label=\alph*.]
\item 0.350 dm$^3$
\item 3.50 dm$^3$
\item 35.0 dm$^3$
\item 350 dm$^3$
\end{enumerate}
\end{minipage} \\
\hline
\vspace{0.2cm}
\textbf{Answer:} 
\[
3.50~\text{L} \times \frac{1~\text{dm}^3}{1~\text{L}} = 3.50~\text{dm}^3
\] \\
\hline
\end{tabular}

\newpage
\noindent
{\large\textbf{Content (fully worked out problems \& answers) continued...}}\\[0.2em]
\vspace{0.2cm}
\begin{tabular}{|p{0.9\linewidth}|}
\hline
\vspace{0.2cm}
\textbf{(65)} A student measures 0.00250 kg of a substance. Express this in grams.\\
\hline
\vspace{0.2cm}
\begin{minipage}[t]{0.85\linewidth}
\begin{enumerate}[label=\alph*.]
\item 0.00250 g
\item 0.0250 g
\item 2.50 g
\item 250 g
\end{enumerate}
\end{minipage}\\
\hline
\vspace{0.2cm}
\textbf{Answer:} 
\[
0.00250~\text{kg} \times \frac{1000~\text{g}}{1~\text{kg}} = 2.50~\text{g}
\] \\ 
\hline
\end{tabular}

\newpage
\noindent
{\large\textbf{Content (fully worked out problems \& answers) continued...}}\\[0.2em]
\vspace{0.2cm}
\begin{tabular}{|p{0.9\linewidth}|}
\hline
\vspace{0.2cm}
\textbf{(66)} A 5.00 mL sample of water weighs 4.98 g. What is its density?\\
\hline
\vspace{0.2cm}
\begin{minipage}[t]{0.85\linewidth}
\begin{enumerate}[label=\alph*.]
\item 0.996 g/mL
\item 1.00 g/mL
\item 1.02 g/mL
\item 1.10 g/mL
\end{enumerate}
\end{minipage}\\
\hline
\vspace{0.2cm}
\textbf{Answer:} 
\[
\rho = \frac{4.98~\text{g}}{5.00~\text{mL}} = 0.996~\text{g/mL}
\] \\ 
\hline
\end{tabular}

\newpage
\noindent
{\large\textbf{Content (fully worked out problems \& answers) continued...}}\\[0.2em]
\vspace{0.2cm}
\begin{tabular}{|p{0.9\linewidth}|}
\hline
\vspace{0.2cm}
\textbf{(67)} Convert steel density 7850 kg/m$^3$ to lb/ft$^3$. (Use \(1~\text{m}^3=35.3147~\text{ft}^3, 1~\text{kg}=2.20462~\text{lb}\))\\
\hline
\vspace{0.2cm}
\begin{minipage}[t]{0.85\linewidth}
\begin{enumerate}[label=\alph*.]
\item 490.0 lb/ft$^3$
\item 785.0 lb/ft$^3$
\item 890 lb/ft$^3$
\item 354.0 lb/ft$^3$
\end{enumerate}
\end{minipage}\\
\hline
\vspace{0.2cm}
\textbf{Answer:} 
\[
7850~\text{kg/m}^3 \times \frac{2.20462~\text{lb}}{1~\text{kg}} \times \frac{1~\text{m}^3}{35.3147~\text{ft}^3} = 490.0~\text{lb/ft}^3
\] \\ 
\hline
\end{tabular}

\newpage
\noindent
{\large\textbf{Content (fully worked out problems \& answers) continued...}}\\[0.2em]
\vspace{0.2cm}
\begin{tabular}{|p{0.9\linewidth}|}
\hline
\vspace{0.2cm}
\textbf{(68)} Multiply $3.445~\text{cm}\times 2.1~\text{cm}$.\
\hline
\vspace{0.2cm}
\begin{minipage}[t]{0.85\linewidth}
\begin{enumerate}[label=\alph*.]
\item 7.2 cm$^2$
\item 7.23 cm$^2$
\item 7.234 cm$^2$
\item 7.235 cm$^2$
\end{enumerate}
\end{minipage}\\
\hline
\vspace{0.2cm}
\textbf{Answer:} $7.2~\text{cm}^2$\\
\[
3.445 \times 2.1 = 7.2345~\text{cm}^2
\]
\[
\text{With correct sig figs (2 sig figs, from 2.1): } 7.2~\text{cm}^2
\] \\
\hline
\end{tabular}
\newpage
\noindent
{\large\textbf{Content (fully worked out problems \& answers) continued...}}\\[0.2em]
\vspace{0.2cm}
\begin{tabular}{|p{0.9\linewidth}|}
\hline
\vspace{0.2cm}
\textbf{(69)} Compute $\dfrac{4.500\times 10^{3}\text{g}}{2.6\text{L}}$. Report with correct sig figs.\\
\hline
\vspace{0.2cm}
\begin{minipage}[t]{0.85\linewidth}
\begin{enumerate}[label=\alph*.]
\item $1.73\times 10^{3}$ g/L
\item $1.7\times 10^{3}$ g/L
\item $1.730\times 10^{3}$ g/L
\item $1.75\times 10^{3}$ g/L
\end{enumerate}
\end{minipage}\\
\hline
\vspace{0.2cm}
\textbf{Answer:} $1.7\times 10^{3}~\text{g/L}$\\
\[
\frac{4.500 \times 10^3}{2.6} = 1730.7692~\text{g/L}
\]

\[
\text{With correct sig figs (2 sig figs, from 2.6 L): } 1.7 \times 10^3~\text{g/L}
\] \\
\hline
\end{tabular}
\newpage
\noindent
{\large\textbf{Content (fully worked out problems \& answers) continued...}}\\[0.2em]
\vspace{0.2cm}
\begin{tabular}{|p{0.9\linewidth}|}
\hline
\vspace{0.2cm}
\textbf{(70)} Add $2.05~\text{g} + 14.1~\text{g} + 3.222~\text{g}$. Report with correct decimal-place precision.\
\hline
\vspace{0.2cm}
\begin{minipage}[t]{0.85\linewidth}
\begin{enumerate}[label=\alph*.]
\item 19.37 g
\item 19.372 g
\item 19.4 g
\item 19 g
\end{enumerate}
\end{minipage}\\
\hline
\vspace{0.2cm}
\textbf{Answer:} $19.372~\text{g}$\\
\[
2.05 + 14.1 + 3.222 = 19.372~\text{g}
\]

\[
\text{Decimal-place precision: 19.372 g (all decimals retained)}
\] \\
\hline
\end{tabular}
\newpage
\noindent
{\large\textbf{Content (fully worked out problems \& answers) continued...}}\\[0.2em]
\vspace{0.2cm}
\begin{tabular}{|p{0.9\linewidth}|}
\hline
\vspace{0.2cm}
\textbf{(71)} Subtract $4.500\times 10^{2}~\text{mL}$ from $1.200\times 10^{3}\text{mL}$. Give the result with correct sig figs.\\
\hline
\vspace{0.2cm}
\begin{minipage}[t]{0.85\linewidth}
\begin{enumerate}[label=\alph*.]
\item $7.50\times 10^{2}$ mL
\item $7.500\times 10^{2}$ mL
\item $7.5\times 10^{2}$ mL
\item $750$ mL
\end{enumerate}
\end{minipage}\\
\hline
\vspace{0.2cm}
\textbf{Answer:} $7.500\times 10^{2}~\text{mL}$\\
\[
1.200 \times 10^3 - 4.500 \times 10^2 = 750~\text{mL}
\]
\[
\text{Correct sig figs (4 sig figs, matching } 1.200 \times 10^3 \text{ and } 4.500 \times 10^2 \text{): } 7.500 \times 10^2~\text{mL}
\] \\

\hline
\end{tabular}
\newpage
\noindent
{\large\textbf{Content (fully worked out problems \& answers) continued...}}\\[0.2em]
\vspace{0.2cm}
\begin{tabular}{|p{0.9\linewidth}|}
\hline
\vspace{0.2cm}
\textbf{(72)} Volume of a cube with side $3.44~\text{cm}$.\\
\hline
\vspace{0.2cm}
\begin{minipage}[t]{0.85\linewidth}
\begin{enumerate}[label=\alph*.]
\item 40.7 cm$^3$
\item 40.71 cm$^3$
\item 41 cm$^3$
\item 40 cm$^3$
\end{enumerate}
\end{minipage}\\
\hline
\vspace{0.2cm}
\textbf{Answer:} $40.7~\text{cm}^3$\\
\[
V = (3.44)^3 = 40.707584~\text{cm}^3
\]

\[
\text{With correct sig figs (3 sig figs from 3.44): } 40.7~\text{cm}^3
\] \\
\hline
\end{tabular}
\newpage
\noindent
{\large\textbf{Content (fully worked out problems \& answers) continued...}}\\[0.2em]
\vspace{0.2cm}
\begin{tabular}{|p{0.9\linewidth}|}
\hline
\vspace{0.2cm}
\textbf{(73)} Multiply $6.230\times 10^{1}\text{cm}\times 4.0\times 10^{-2}\text{cm}$. Give result with correct sig figs.\\
\hline
\vspace{0.2cm}
\begin{minipage}[t]{0.85\linewidth}
\begin{enumerate}[label=\alph*.]
\item 2.492 cm$^2$
\item 2.49 cm$^2$
\item 2.5 cm$^2$
\item 2.50 cm$^2$
\end{enumerate}
\end{minipage}\\
\hline
\vspace{0.2cm}
\textbf{Answer:} $2.5~\text{cm}^2$\\
\[
6.230 \times 10^1 \times 4.0 \times 10^{-2} = 2.492~\text{cm}^2
\]

\[
\text{Correct sig figs (2 sig figs, from 4.0): } 2.5~\text{cm}^2
\] \\
\hline
\end{tabular}
\newpage
\noindent
{\large\textbf{Content (fully worked out problems \& answers) continued...}}\\[0.2em]
\vspace{0.2cm}
\begin{tabular}{|p{0.9\linewidth}|}
\hline
\vspace{0.2cm}
\textbf{(74)} A mass of $12.35~\text{g}$ and a mass of $1.2~\text{g}$ are combined. What is the combined mass with correct sig figs?\\
\hline
\vspace{0.2cm}
\begin{minipage}[t]{0.85\linewidth}
\begin{enumerate}[label=\alph*.]
\item 13.55 g
\item 13.6 g
\item 13.5 g
\item 14 g
\end{enumerate}
\end{minipage}\\
\hline
\vspace{0.2cm}
\textbf{Answer:} $13.55~\text{g}$\\
\[
12.35 + 1.2 = 13.55~\text{g}
\]

\[
\text{Decimal precision determined by 1.2 (2 sig figs): } 13.55~\text{g}
\] \\
\hline
\end{tabular}
\newpage
\noindent
{\large\textbf{Content (fully worked out problems \& answers) continued...}}\\[0.2em]
\vspace{0.2cm}
\begin{tabular}{|p{0.9\linewidth}|}
\hline
\vspace{0.2cm}
\textbf{(75)} A length $4.80~\text{km}$ has an exact addition of $250~\text{m}$. Express the total in km with correct sig figs.\\
\hline
\vspace{0.2cm}
\begin{minipage}[t]{0.85\linewidth}
\begin{enumerate}[label=\alph*.]
\item 5.05 km
\item 5.050 km
\item 5.1 km
\item 5.0 km
\end{enumerate}
\end{minipage}\\
\hline
\vspace{0.2cm}
\textbf{Answer:} $5.05~\text{km}$\\
\[
250~\text{m} = 0.250~\text{km}
\]

\[
4.80 + 0.250 = 5.05~\text{km}
\] \\
\hline
\end{tabular}
\newpage
\noindent
{\large\textbf{Content (fully worked out problems \& answers) continued...}}\\[0.2em]
\vspace{0.2cm}
\begin{tabular}{|p{0.9\linewidth}|}
\hline
\vspace{0.2cm}
\textbf{(76)} Convert $9.60\times 10^{1}\text{mph}$ to m/s and report with correct sig figs.\\
\hline
\vspace{0.2cm}
\begin{minipage}[t]{0.85\linewidth}
\begin{enumerate}[label=\alph*.]
\item 42.9 m/s
\item 43.0 m/s
\item 42.94 m/s
\item 42 m/s
\end{enumerate}
\end{minipage}\\
\hline
\vspace{0.2cm}
\textbf{Answer:} $42.9~\text{m/s}$\\
\[
1~\text{mph} = 0.44704~\text{m/s}
\]

\[
96.0 \times 0.44704 = 42.892~\text{m/s}
\]

\[
\text{With 3 sig figs from 9.60×10¹: } 42.9~\text{m/s}
\] \\
\hline
\end{tabular}
\newpage
\noindent
{\large\textbf{Content (fully worked out problems \& answers) continued...}}\\[0.2em]
\vspace{0.2cm}
\begin{tabular}{|p{0.9\linewidth}|}
\hline
\vspace{0.2cm}
\textbf{(77)} A beaker has $125.0~\text{mL}$; an exact $10.00~\text{mL}$ pipet adds more. What is the new volume with correct sig figs?\\
\hline
\vspace{0.2cm}
\begin{minipage}[t]{0.85\linewidth}
\begin{enumerate}[label=\alph*.]
\item 135 mL
\item 135.0 mL
\item 135.00 mL
\item 135.000 mL
\end{enumerate}
\end{minipage}\\
\hline
\vspace{0.2cm}
\textbf{Answer:} $135.0~\text{mL}$\\
\[
125.0 + 10.00 = 135.0~\text{mL}
\]

\[
\text{Correct sig figs: 135.0 mL (4 sig figs total)}
\] \\
\hline
\end{tabular}
\newpage
\noindent
{\large\textbf{Content (fully worked out problems \& answers) continued...}}\\[0.2em]
\vspace{0.2cm}
\begin{tabular}{|p{0.9\linewidth}|}
\hline
\vspace{0.2cm}
\textbf{(78)} Convert $3.25~\text{lb}$ to grams using $1~\text{lb}=453.6~\text{g}$.\\
\hline
\vspace{0.2cm}
\begin{minipage}[t]{0.85\linewidth}
\begin{enumerate}[label=\alph*.]
\item 1474 g
\item $1.47 \times 10^{3}$ g
\item $1.474 \times 10^{3}$ g
\item $1.48 \times 10^{3}$ g
\end{enumerate}
\end{minipage}\\
\hline
\vspace{0.2cm}
\textbf{Answer:} $1.47\times 10^{3}~\text{g}$\\

\[
3.25 \times 453.6 = 1474.2~\text{g}
\]

\[
\text{Correct sig figs (3 sig figs from 3.25): } 1.47 \times 10^3~\text{g}
\] \\
\hline
\end{tabular}
\newpage
\noindent
{\large\textbf{Content (fully worked out problems \& answers) continued...}}\\[0.2em]
\vspace{0.2cm}
\begin{tabular}{|p{0.9\linewidth}|}
\hline
\vspace{0.2cm}
\textbf{(79)} Express $7.400\times 10^{-2}\text{L}$ in cm$^3$ (1 L = 1000 cm$^3$).\\
\hline
\vspace{0.2cm}
\begin{minipage}[t]{0.85\linewidth}
\begin{enumerate}[label=\alph*.]
\item 74 cm$^3$
\item 74.0 cm$^3$
\item 74.00 cm$^3$
\item 74.000 cm$^3$
\end{enumerate}
\end{minipage}\\
\hline
\vspace{0.2cm}
\textbf{Answer:} $74.00~\text{cm}^3$\\
\[
7.400 \times 10^{-2} \times 1000 = 74.0~\text{cm}^3
\]

\[
\text{Correct sig figs (4 sig figs from 7.400×10⁻²): } 74.00~\text{cm}^3
\] \\
\hline
\end{tabular}
\newpage
\noindent
{\large\textbf{Content (fully worked out problems \& answers) continued...}}\\[0.2em]
\vspace{0.2cm}
\begin{tabular}{|p{0.9\linewidth}|}
\hline
\vspace{0.2cm}
\textbf{(80)} You have $3.450\times10^{2}~\text{mL}$ of solution and add $1.2\times10^{3}~\text{mL}$ solvent. Total volume? \\
\hline
\vspace{0.2cm}
\begin{minipage}[t]{0.85\linewidth}
\begin{enumerate}[label=\alph*.]
\item 1545 mL
\item $1.55 \times 10^{3}$ mL
\item $1.545 \times 10^{3}$ mL
\item $1.540 \times 10^{3}$ mL
\end{enumerate}

\end{minipage}\\
\hline
\vspace{0.2cm}
\textbf{Answer:} $1.545\times10^{3}~\text{mL}$\\
\[
345.0 + 1200 = 1545.0~\text{mL}
\]

\[
\text{Correct sig figs: } 1.545 \times 10^3~\text{mL}
\] \\
\hline
\end{tabular}
\newpage
\noindent
{\large\textbf{Content (fully worked out problems \& answers) continued...}}\\[0.2em]
\vspace{0.2cm}
\begin{tabular}{|p{0.9\linewidth}|}
\hline
\vspace{0.2cm}
\textbf{(81)} A bottle contains $1.003\times10^{3}\text{g}$ of NaCl. Remove $125.0\text{g}$. Mass remaining? \\
\hline
\vspace{0.2cm}
\begin{minipage}[t]{0.85\linewidth}
\begin{enumerate}[label=\alph*.]
\item 878 g
\item 878.0 g
\item 877.95 g
\item 880 g
\end{enumerate}
\end{minipage}\\
\hline
\vspace{0.2cm}
\textbf{Answer:} $878.0~\text{g}$\\
\[
1003 - 125.0 = 878.0~\text{g}
\]

\[
\text{Correct sig figs: } 878.0~\text{g}
\] \\
\hline
\end{tabular}
\newpage
\noindent
{\large\textbf{Content (fully worked out problems \& answers) continued...}}\\[0.2em]
\vspace{0.2cm}
\begin{tabular}{|p{0.9\linewidth}|}
\hline
\vspace{0.2cm}
\textbf{(82)} Two glass rods: $12.45~\text{cm}$ and $8.0~\text{cm}$. Combined length? \\
\hline
\vspace{0.2cm}
\begin{minipage}[t]{0.85\linewidth}
\begin{enumerate}[label=\alph*.]
\item 20.45 cm
\item 20.5 cm
\item 20.4 cm
\item 21 cm
\end{enumerate}
\end{minipage}\\
\hline
\vspace{0.2cm}
\textbf{Answer:} $20.45~\text{cm}$\\
\[
12.45 + 8.0 = 20.45~\text{cm}
\]

\[
\text{Correct sig figs (decimal place dictated by 8.0): } 20.45~\text{cm}
\] \\
\hline
\end{tabular}
\newpage
\noindent
{\large\textbf{Content (fully worked out problems \& answers) continued...}}\\[0.2em]
\vspace{0.2cm}
\begin{tabular}{|p{0.9\linewidth}|}
\hline
\vspace{0.2cm}
\textbf{(83)} Stock solution volume $2.500\times10^{3}\text{mL}$. Add 45 mL water. Final volume?\\
\hline
\vspace{0.2cm}
\begin{minipage}[t]{0.85\linewidth}
\begin{enumerate}[label=\alph*.]
\item 2545.0 mL
\item $2.545 \times 10^{3}$ mL
\item $2.55 \times 10^{3}$ mL
\item $2.54 \times 10^{3}$ mL
\end{enumerate}

\end{minipage}\\
\hline
\vspace{0.2cm}
\textbf{Answer:} $2.545\times10^{3}~\text{mL}$\\
\[
2500 + 45 = 2545~\text{mL}
\]

\[
\text{Correct sig figs: } 2.545 \times 10^3~\text{mL}
\] \\
\hline
\end{tabular}
\newpage
\noindent
{\large\textbf{Content (fully worked out problems \& answers) continued...}}\\[0.2em]
\vspace{0.2cm}
\begin{tabular}{|p{0.9\linewidth}|}
\hline
\vspace{0.2cm}
\textbf{(84)} Buret reading drops from $12.004~\text{mL}$ to $11.9965~\text{mL}$. Subtract $0.0045~\text{mL}$ from 12.004. Report correct sig figs.\
\hline
\vspace{0.2cm}
\begin{minipage}[t]{0.85\linewidth}
\begin{enumerate}[label=\alph*.]
\item 11.9995 mL
\item 12.000 mL
\item 12.00 mL
\item 12.0 mL
\end{enumerate}
\end{minipage}\\
\hline
\vspace{0.2cm}
\textbf{Answer:} $11.9995~\text{mL}$\\
\[
12.004 - 0.0045 = 11.9995~\text{mL}
\]

\[
\text{Correct sig figs (all digits meaningful): } 11.9995~\text{mL}
\] \\
\hline
\end{tabular}
\newpage

\noindent
{\large\textbf{Content (fully worked out problems \& answers) continued...}}\\[0.2em]
\begin{tabular}{|p{0.9\linewidth}|}
\hline
\vspace{0.2cm}
\textbf{(85)} A metal foil width is $2.340~\text{mm}$; coating adds $0.002~\text{mm}$. Combined width?\\
\hline
\vspace{0.2cm}
\begin{minipage}[t]{0.85\linewidth}
\begin{enumerate}[label=\alph*.]
\item 2.342 mm
\item 2.34 mm
\item 2.344 mm
\item 2.345 mm
\end{enumerate}
\end{minipage}\\
\hline
\vspace{0.2cm}
\textbf{Answer:} $2.342~\text{mm}$\\
\[
2.340 + 0.002 = 2.342~\text{mm}
\]

\[
\text{Correct sig figs: } 2.342~\text{mm}
\] \\
\hline
\end{tabular}
\newpage
\noindent
{\large\textbf{Content (fully worked out problems \& answers) continued...}}\\[0.2em]
\begin{tabular}{|p{0.9\linewidth}|}
\hline
\vspace{0.2cm}
\textbf{(86)} Wet sample mass $34.55~\text{g}$; dry $33.1~\text{g}$. Mass lost?\\
\hline
\vspace{0.2cm}
\begin{minipage}[t]{0.85\linewidth}
\begin{enumerate}[label=\alph*.]
\item 1.45 g
\item 1.5 g
\item 1.50 g
\item 1.4 g
\end{enumerate}
\end{minipage}\\
\hline
\vspace{0.2cm}
\textbf{Answer:} $1.45~\text{g}$\\
\[
34.55 - 33.1 = 1.45~\text{g}
\]

\[
\text{Correct sig figs: } 1.45~\text{g}
\] \\
\hline
\end{tabular}
\newpage
\noindent
{\large\textbf{Content (fully worked out problems \& answers) continued...}}\\[0.2em]
\begin{tabular}{|p{0.9\linewidth}|}
\hline
\vspace{0.2cm}
\textbf{(87)} Reaction releases $4.50\times10^{4}~\text{J}$ then $1.2\times10^{3}~\text{J}$. Total heat?\\
\hline
\vspace{0.2cm}
\begin{minipage}[t]{0.85\linewidth}
\begin{enumerate}[label=\alph*.]
\item $4.62 \times 10^{4}$ J
\item $4.623 \times 10^{4}$ J
\item $4.6 \times 10^{4}$ J
\item $4.625 \times 10^{4}$ J
\end{enumerate}
\end{minipage}\\
\hline
\vspace{0.2cm}
\textbf{Answer:} $4.62\times10^{4}~\text{J}$\\
\[
4.50 \times 10^4 + 1.2 \times 10^3 = 46200~\text{J} = 4.62 \times 10^4~\text{J}
\] \\
\hline
\end{tabular}
\newpage
\noindent
{\large\textbf{Content (fully worked out problems \& answers) continued...}}\\[0.2em]
\begin{tabular}{|p{0.9\linewidth}|}
\hline
\vspace{0.2cm}
\textbf{(88)} Subtract $6.020\times10^{23}$ molecules from $6.022\times10^{23}$. Difference?\\
\hline
\vspace{0.2cm}
\begin{minipage}[t]{0.85\linewidth}
\begin{enumerate}[label=\alph*.]
\item $2.0 \times 10^{20}$ molecules
\item $2 \times 10^{20}$ molecules
\item $0.002 \times 10^{23}$ molecules
\item $2.00 \times 10^{20}$ molecules
\end{enumerate}
\end{minipage}\\
\hline
\vspace{0.2cm}
\textbf{Answer:} $2.0\times10^{20}$ molecules\\
\[
6.022 \times 10^{23} - 6.020 \times 10^{23} = 0.002 \times 10^{23} = 2.0 \times 10^{20}~\text{molecules}
\] \\
\hline
\end{tabular}
\newpage
\noindent
{\large\textbf{Content (fully worked out problems \& answers) continued...}}\\[0.2em]
\begin{tabular}{|p{0.9\linewidth}|}
\hline
\vspace{0.2cm}
\textbf{(89)} Lab distances: $12.3~\text{m}$ and $8.45~\text{m}$. Total?\\
\hline
\vspace{0.2cm}
\begin{minipage}[t]{0.85\linewidth}
\begin{enumerate}[label=\alph*.]
\item 20.8 m
\item 20.75 m
\item 20.7 m
\item 21 m
\end{enumerate}
\end{minipage}\\
\hline
\vspace{0.2cm}
\textbf{Answer:} $20.75~\text{m}$\\
\[
12.3 + 8.45 = 20.75~\text{m}
\] \\
\hline
\end{tabular}
\newpage
\noindent
{\large\textbf{Content (fully worked out problems \& answers) continued...}}\\[0.2em]
\begin{tabular}{|p{0.9\linewidth}|}
\hline
\vspace{0.2cm}
\textbf{(90)} Add $0.00845~\text{m}$ and $0.045~\text{m}$ for a capillary. Sum?\\
\hline
\vspace{0.2cm}
\begin{minipage}[t]{0.85\linewidth}
\begin{enumerate}[label=\alph*.]
\item 0.053 m
\item 0.05345 m
\item 0.0535 m
\item 0.054 m
\end{enumerate}
\end{minipage}\\
\hline
\vspace{0.2cm}
\textbf{Answer:} $0.05345~\text{m}$\\
\[
0.00845 + 0.045 = 0.05345~\text{m}
\] \\ 
\hline
\end{tabular}
\newpage
\noindent
{\large\textbf{Content (fully worked out problems \& answers) continued...}}\\[0.2em]
\begin{tabular}{|p{0.9\linewidth}|}
\hline
\vspace{0.2cm}
\textbf{(91)} You have $1.0000~\text{L}$ solvent; remove $1.00~\text{L}$. Remaining volume?\\
\hline
\vspace{0.2cm}
\begin{minipage}[t]{0.85\linewidth}
\begin{enumerate}[label=\alph*.]
\item 0.0000 L
\item 0.00 L
\item 0 L
\item 0.0 L
\end{enumerate}
\end{minipage}\\
\hline
\vspace{0.2cm}
\textbf{Answer:} $0.00~\text{L}$\\
\[
1.0000 - 1.00 = 0.00~\text{L}
\] \\ 
\hline
\end{tabular}
\newpage
\noindent
{\large\textbf{Content (fully worked out problems \& answers) continued...}}\\[0.2em]
\begin{tabular}{|p{0.9\linewidth}|}
\hline
\vspace{0.2cm}
\textbf{(92)} Combine metal rods: $0.452~\text{m}$ and $0.0035~\text{m}$. Total length?\\
\hline
\vspace{0.2cm}
\begin{minipage}[t]{0.85\linewidth}
\begin{enumerate}[label=\alph*.]
\item 0.456 m
\item 0.4555 m
\item 0.455 m
\item 0.457 m
\end{enumerate}
\end{minipage}\\
\hline
\vspace{0.2cm}
\textbf{Answer:} $0.4587~\text{m}$\\
\[
0.452 + 0.0035 = 0.4555~\text{m}
\]

\[
\text{Correct sig figs (4 sig figs from 0.0035): } 0.4555~\text{m}
\] \\ 
\hline
\end{tabular}
\newpage
\noindent
{\large\textbf{Content (fully worked out problems& answers) continued...}}\\[0.2em]
\vspace{0.2cm}
\begin{tabular}{|p{0.9\linewidth}|}
\hline
\vspace{0.2cm}
\textbf{(93)} Which state of matter has a definite shape and volume?\\
\begin{minipage}[t]{0.85\linewidth}
\begin{enumerate}[label=\alph*.]
\item Solid
\item Liquid
\item Gas
\item Plasma
\end{enumerate}
\end{minipage}\\
\hline
\vspace{0.2cm}
\textbf{Answer:} a. Solid\\
\hline
\end{tabular}

\newpage
\noindent
{\large\textbf{Content (fully worked out problems& answers) continued...}}\\[0.2em]
\vspace{0.2cm}
\begin{tabular}{|p{0.9\linewidth}|}
\hline
\vspace{0.2cm}
\textbf{(94)} What is a chemical property?\\
\begin{minipage}[t]{0.85\linewidth}
\begin{enumerate}[label=\alph*.]
\item A characteristic observed without changing the substance
\item A characteristic that describes how a substance reacts with other substances
\item The color of a substance
\item The mass of a substance
\end{enumerate}
\end{minipage}\\
\hline
\vspace{0.2cm}
\textbf{Answer:} b. A characteristic that describes how a substance reacts with other substances\\
\hline
\end{tabular}

\newpage
\noindent
{\large\textbf{Content (fully worked out problems& answers) continued...}}\\[0.2em]
\vspace{0.2cm}
\begin{tabular}{|p{0.9\linewidth}|}
\hline
\vspace{0.2cm}
\textbf{(95)} Which of the following is a physical change?\\
\begin{minipage}[t]{0.85\linewidth}
\begin{enumerate}[label=\alph*.]
\item Rusting of iron
\item Burning wood
\item Melting ice
\item Baking a cake
\end{enumerate}
\end{minipage}\\
\hline
\vspace{0.2cm}
\textbf{Answer:} c. Melting ice\\
\hline
\end{tabular}

\newpage
\noindent
{\large\textbf{Content (fully worked out problems& answers) continued...}}\\[0.2em]
\vspace{0.2cm}
\begin{tabular}{|p{0.9\linewidth}|}
\hline
\vspace{0.2cm}
\textbf{(96)} Which particle is found in the nucleus of an atom?\\
\begin{minipage}[t]{0.85\linewidth}
\begin{enumerate}[label=\alph*.]
\item Electron
\item Proton
\item Neutron
\item Both b and c
\end{enumerate}
\end{minipage}\\
\hline
\vspace{0.2cm}
\textbf{Answer:} d. Both b and c\\
\hline
\end{tabular}

\newpage
\noindent
{\large\textbf{Content (fully worked out problems& answers) continued...}}\\[0.2em]
\vspace{0.2cm}
\begin{tabular}{|p{0.9\linewidth}|}
\hline
\vspace{0.2cm}
\textbf{(97)} What is an element?\\
\begin{minipage}[t]{0.85\linewidth}
\begin{enumerate}[label=\alph*.]
\item A substance made of two or more different atoms chemically bonded
\item A pure substance made of only one type of atom
\item A mixture of compounds
\item A solution
\end{enumerate}
\end{minipage}\\
\hline
\vspace{0.2cm}
\textbf{Answer:} b. A pure substance made of only one type of atom\\
\hline
\end{tabular}

\newpage
\noindent
{\large\textbf{Content (fully worked out problems& answers) continued...}}\\[0.2em]
\vspace{0.2cm}
\begin{tabular}{|p{0.9\linewidth}|}
\hline
\vspace{0.2cm}
\textbf{(98)} Which of the following is an example of a chemical change?\\
\begin{minipage}[t]{0.85\linewidth}
\begin{enumerate}[label=\alph*.]
\item Boiling water
\item Dissolving sugar in water
\item Burning paper
\item Freezing juice
\end{enumerate}
\end{minipage}\\
\hline
\vspace{0.2cm}
\textbf{Answer:} c. Burning paper\\
\hline
\end{tabular}

\newpage
\noindent
{\large\textbf{Content (fully worked out problems& answers) continued...}}\\[0.2em]
\vspace{0.2cm}
\begin{tabular}{|p{0.9\linewidth}|}
\hline
\vspace{0.2cm}
\textbf{(99)} Which phase of matter has particles that move freely and fill the shape of their container?\\
\begin{minipage}[t]{0.85\linewidth}
\begin{enumerate}[label=\alph*.]
\item Solid
\item Liquid
\item Gas
\item Plasma
\end{enumerate}
\end{minipage}\\
\hline
\vspace{0.2cm}
\textbf{Answer:} c. Gas\\
\hline
\end{tabular}

\newpage
\noindent
{\large\textbf{Content (fully worked out problems& answers) continued...}}\\[0.2em]
\vspace{0.2cm}
\begin{tabular}{|p{0.9\linewidth}|}
\hline
\vspace{0.2cm}
\textbf{(100)} Which of the following is a physical property?\\
\begin{minipage}[t]{0.85\linewidth}
\begin{enumerate}[label=\alph*.]
\item Flammability
\item Reactivity with acid
\item Density
\item Oxidation
\end{enumerate}
\end{minipage}\\
\hline
\vspace{0.2cm}
\textbf{Answer:} c. Density\\
\hline
\end{tabular}
\newpage
\newpage
\noindent
{\large\textbf{Content (fully worked out problems& answers) continued...}}\\[0.2em]
\vspace{0.2cm}
\begin{tabular}{|p{0.9\linewidth}|}
\hline
\vspace{0.2cm}
\textbf{(101)} Which state of matter takes both the shape and volume of its container?\\
\begin{minipage}[t]{0.85\linewidth}
\begin{enumerate}[label=\alph*.]
\item Solid
\item Liquid
\item Gas
\item Plasma
\end{enumerate}
\end{minipage}\\
\hline
\vspace{0.2cm}
\textbf{Answer:} b. Liquid\\
\hline
\end{tabular}

\newpage
\noindent
{\large\textbf{Content (fully worked out problems& answers) continued...}}\\[0.2em]
\vspace{0.2cm}
\begin{tabular}{|p{0.9\linewidth}|}
\hline
\vspace{0.2cm}
\textbf{(102)} What is the smallest particle of an element that retains its chemical properties?\\
\begin{minipage}[t]{0.85\linewidth}
\begin{enumerate}[label=\alph*.]
\item Molecule
\item Atom
\item Ion
\item Compound
\end{enumerate}
\end{minipage}\\
\hline
\vspace{0.2cm}
\textbf{Answer:} b. Atom\\
\hline
\end{tabular}

\newpage
\noindent
{\large\textbf{Content (fully worked out problems& answers) continued...}}\\[0.2em]
\vspace{0.2cm}
\begin{tabular}{|p{0.9\linewidth}|}
\hline
\vspace{0.2cm}
\textbf{(103)} Which of the following is a chemical property?\\
\begin{minipage}[t]{0.85\linewidth}
\begin{enumerate}[label=\alph*.]
\item Boiling point
\item Color
\item Reactivity with water
\item Density
\end{enumerate}
\end{minipage}\\
\hline
\vspace{0.2cm}
\textbf{Answer:} c. Reactivity with water\\
\hline
\end{tabular}

\newpage
\noindent
{\large\textbf{Content (fully worked out problems& answers) continued...}}\\[0.2em]
\vspace{0.2cm}
\begin{tabular}{|p{0.9\linewidth}|}
\hline
\vspace{0.2cm}
\textbf{(104)} Which of the following is an example of a physical change?\\
\begin{minipage}[t]{0.85\linewidth}
\begin{enumerate}[label=\alph*.]
\item Tarnishing of silver
\item Cutting paper
\item Burning coal
\item Rusting iron
\end{enumerate}
\end{minipage}\\
\hline
\vspace{0.2cm}
\textbf{Answer:} b. Cutting paper\\
\hline
\end{tabular}

\newpage
\noindent
{\large\textbf{Content (fully worked out problems& answers) continued...}}\\[0.2em]
\vspace{0.2cm}
\begin{tabular}{|p{0.9\linewidth}|}
\hline
\vspace{0.2cm}
\textbf{(105)} What is the charge of an electron?\\
\begin{minipage}[t]{0.85\linewidth}
\begin{enumerate}[label=\alph*.]
\item Positive
\item Negative
\item Neutral
\item Variable
\end{enumerate}
\end{minipage}\\
\hline
\vspace{0.2cm}
\textbf{Answer:} b. Negative\\
\hline
\end{tabular}

\newpage
\noindent
{\large\textbf{Content (fully worked out problems& answers) continued...}}\\[0.2em]
\vspace{0.2cm}
\begin{tabular}{|p{0.9\linewidth}|}
\hline
\vspace{0.2cm}
\textbf{(106)} Which of the following is NOT a phase of matter?\\
\begin{minipage}[t]{0.85\linewidth}
\begin{enumerate}[label=\alph*.]
\item Solid
\item Liquid
\item Gas
\item Solution
\end{enumerate}
\end{minipage}\\
\hline
\vspace{0.2cm}
\textbf{Answer:} d. Solution\\
\hline
\end{tabular}

\newpage
\noindent
{\large\textbf{Content (fully worked out problems& answers) continued...}}\\[0.2em]
\vspace{0.2cm}
\begin{tabular}{|p{0.9\linewidth}|}
\hline
\vspace{0.2cm}
\textbf{(107)} What is the center of an atom called?\\
\begin{minipage}[t]{0.85\linewidth}
\begin{enumerate}[label=\alph*.]
\item Electron cloud
\item Nucleus
\item Proton
\item Neutron
\end{enumerate}
\end{minipage}\\
\hline
\vspace{0.2cm}
\textbf{Answer:} b. Nucleus\\
\hline
\end{tabular}

\newpage
\noindent
{\large\textbf{Content (fully worked out problems& answers) continued...}}\\[0.2em]
\vspace{0.2cm}
\begin{tabular}{|p{0.9\linewidth}|}
\hline
\vspace{0.2cm}
\textbf{(108)} Which of the following is an element symbol?\\
\begin{minipage}[t]{0.85\linewidth}
\begin{enumerate}[label=\alph*.]
\item H₂O
\item CO₂
\item O
\item NaCl
\end{enumerate}
\end{minipage}\\
\hline
\vspace{0.2cm}
\textbf{Answer:} c. O\\
\hline
\end{tabular}

\newpage
\noindent
{\large\textbf{Content (fully worked out problems& answers) continued...}}\\[0.2em]
\vspace{0.2cm}
\begin{tabular}{|p{0.9\linewidth}|}
\hline
\vspace{0.2cm}
\textbf{(109)} Which of the following is an example of a chemical change?\\
\begin{minipage}[t]{0.85\linewidth}
\begin{enumerate}[label=\alph*.]
\item Freezing water
\item Dissolving salt in water
\item Cooking an egg
\item Breaking a glass
\end{enumerate}
\end{minipage}\\
\hline
\vspace{0.2cm}
\textbf{Answer:} c. Cooking an egg\\
\hline
\end{tabular}

\newpage
\noindent
{\large\textbf{Content (fully worked out problems& answers) continued...}}\\[0.2em]
\vspace{0.2cm}
\begin{tabular}{|p{0.9\linewidth}|}
\hline
\vspace{0.2cm}
\textbf{(110)} Which particle determines the chemical identity of an element?\\
\begin{minipage}[t]{0.85\linewidth}
\begin{enumerate}[label=\alph*.]
\item Electron
\item Neutron
\item Proton
\item All of the above
\end{enumerate}
\end{minipage}\\
\hline
\vspace{0.2cm}
\textbf{Answer:} c. Proton\\
\hline
\end{tabular}
\newpage
\newpage
\noindent
{\large\textbf{Content (fully worked out problems& answers) continued...}}\\[0.2em]
\vspace{0.2cm}
\begin{tabular}{|p{0.9\linewidth}|}
\hline
\vspace{0.2cm}
\textbf{(111)} What is the main difference between a physical change and a chemical change?\\
\begin{minipage}[t]{0.85\linewidth}
\begin{enumerate}[label=\alph*.]
\item Physical changes alter the composition; chemical changes do not
\item Physical changes do not alter the composition; chemical changes do
\item Both alter the composition
\item Neither alters the composition
\end{enumerate}
\end{minipage}\\
\hline
\vspace{0.2cm}
\textbf{Answer:} b. Physical changes do not alter the composition; chemical changes do\\
\hline
\end{tabular}

\newpage
\noindent
{\large\textbf{Content (fully worked out problems& answers) continued...}}\\[0.2em]
\vspace{0.2cm}
\begin{tabular}{|p{0.9\linewidth}|}
\hline
\vspace{0.2cm}
\textbf{(112)} Which of the following is an example of a gas at room temperature?\\
\begin{minipage}[t]{0.85\linewidth}
\begin{enumerate}[label=\alph*.]
\item Oxygen
\item Iron
\item Water
\item Mercury
\end{enumerate}
\end{minipage}\\
\hline
\vspace{0.2cm}
\textbf{Answer:} a. Oxygen\\
\hline
\end{tabular}

\newpage
\noindent
{\large\textbf{Content (fully worked out problems& answers) continued...}}\\[0.2em]
\vspace{0.2cm}
\begin{tabular}{|p{0.9\linewidth}|}
\hline
\vspace{0.2cm}
\textbf{(113)} What is the basic unit of matter?\\
\begin{minipage}[t]{0.85\linewidth}
\begin{enumerate}[label=\alph*.]
\item Atom
\item Molecule
\item Element
\item Compound
\end{enumerate}
\end{minipage}\\
\hline
\vspace{0.2cm}
\textbf{Answer:} a. Atom\\
\hline
\end{tabular}

\newpage
\noindent
{\large\textbf{Content (fully worked out problems& answers) continued...}}\\[0.2em]
\vspace{0.2cm}
\begin{tabular}{|p{0.9\linewidth}|}
\hline
\vspace{0.2cm}
\textbf{(114)} Which of the following is a mixture?\\
\begin{minipage}[t]{0.85\linewidth}
\begin{enumerate}[label=\alph*.]
\item Salt (NaCl)
\item Sugar (C₆H₁₂O₆)
\item Air
\item Oxygen
\end{enumerate}
\end{minipage}\\
\hline
\vspace{0.2cm}
\textbf{Answer:} c. Air\\
\hline
\end{tabular}

\newpage
\noindent
{\large\textbf{Content (fully worked out problems& answers) continued...}}\\[0.2em]
\vspace{0.2cm}
\begin{tabular}{|p{0.9\linewidth}|}
\hline
\vspace{0.2cm}
\textbf{(115)} Which of the following describes an atom?\\
\begin{minipage}[t]{0.85\linewidth}
\begin{enumerate}[label=\alph*.]
\item Smallest particle of a compound
\item Smallest particle of an element that retains its properties
\item A molecule made of different elements
\item A type of mixture
\end{enumerate}
\end{minipage}\\
\hline
\vspace{0.2cm}
\textbf{Answer:} b. Smallest particle of an element that retains its properties\\
\hline
\end{tabular}

\newpage
\noindent
{\large\textbf{Content (fully worked out problems& answers) continued...}}\\[0.2em]
\vspace{0.2cm}
\begin{tabular}{|p{0.9\linewidth}|}
\hline
\vspace{0.2cm}
\textbf{(116)} Which of the following is a property that can be observed without changing the chemical composition?\\
\begin{minipage}[t]{0.85\linewidth}
\begin{enumerate}[label=\alph*.]
\item Flammability
\item Reactivity
\item Melting point
\item Oxidation
\end{enumerate}
\end{minipage}\\
\hline
\vspace{0.2cm}
\textbf{Answer:} c. Melting point\\
\hline
\end{tabular}

\newpage
\noindent
{\large\textbf{Content (fully worked out problems& answers) continued...}}\\[0.2em]
\vspace{0.2cm}
\begin{tabular}{|p{0.9\linewidth}|}
\hline
\vspace{0.2cm}
\textbf{(117)} Which phase of matter has a definite volume but no definite shape?\\
\begin{minipage}[t]{0.85\linewidth}
\begin{enumerate}[label=\alph*.]
\item Solid
\item Liquid
\item Gas
\item Plasma
\end{enumerate}
\end{minipage}\\
\hline
\vspace{0.2cm}
\textbf{Answer:} b. Liquid\\
\hline
\end{tabular}
\newpage

\section*{\underline{Activity 1}}
\begin{tabular}{|p{4.2cm}|p{9.8cm}|}
\hline
\textbf{Collaborative Learning Techniques} & Assigned Discussion Leader\\
\hline
\textbf{Learning Strategy} & Family Feud \\
\hline
\textbf{Time} & 12:10---12:30\\
\hline
\textbf{Materials} & Calculator, writing utensil, paper / tablet, conversion sheets \\
\hline
\textbf{Lecture/Textbook/Old Plans/Slide References} & Chapter 1: Basics of Matter, Chapter 2: Measurements\\
\hline
\end{tabular}

\vspace{0.5cm}

{\large\textbf{Definition}}\\
\noindent
\begin{tabular}{|p{0.9\linewidth}|}
\hline
\begin{enumerate}
    \item\textbf{Assigned Discussion Leader:} One person in the group is asked to present on a topic or review material for the group and then lead the discussion for the group. This person should not be the regular group leader.
    \item\textbf{Family Feud:} This will work well as a test or quiz review for concepts and making connections between terms. Have students write down one term that comes to mind when a broad category is given. Create a list of the top 4 or 5 and use those as "answers" for the game. During the session, dividie studetns into two groups. Given the topic (broad category from the survey), students can decide wheher they want to play or pass. If they play, they alternate turns calling out a term they think is relevant. If they are correct, they get the points. If not, the next student attempts. Continue in this fashion until all terms are guessed and explained or three incorrect guesses are made. If three incorrect guesses occur, the other team tgets the chance to finish the topic.  
    \end{enumerate}\\
\hline
\end{tabular}
\vspace{0.2cm}
\newpage
{\large\textbf{Instructions (please provide a\hl{DETAILED} activity breakdown below (and/or attached}}\\
\noindent
\begin{tabular}{|p{0.9\linewidth}|}
\hline
\vspace{0.2cm}
\begin{enumerate}
    \item Students are organized into teams and reminded that the game covers significant figures, unit conversions, basics of matter, story problems, and fun chemistry facts.
    \item Each team takes turns selecting a category and point value. The team has a set time (e.g., 30–60 seconds) to discuss and provide an answer.
    \item The SI leader monitors discussions, ensuring all students participate and encouraging teams to explain their reasoning before answering.
    \item If a team answers incorrectly, other teams have the opportunity to discuss and “steal” the question.
    \item Points are awarded for correct answers and tracked on the board or a visible scoreboard, helping students see progress and encouraging friendly competition.
    \item Between rounds, the SI leader pauses to clarify any misconceptions, review key concepts, and provide real-life examples to connect the question to practical applications.
    \item After all questions are answered, the SI leader leads a short debrief, asking teams to reflect on what concepts were challenging, what strategies helped, and how preparation and teamwork influenced performance.
\end{enumerate} \\

\hline
\end{tabular}\\
\vspace{0.5cm}
\noindent
\newpage
{\large\textbf{Checking for Understanding (questions \& answers)} \\
\begin{tabular}{|p{0.9\linewidth}|}
\hline
\paragraph{Using a Jeopardy Game to Review Chemistry Concepts}
Engaging students in a Jeopardy-style game serves as an effective and interactive method for reviewing essential chemistry topics while simultaneously assessing comprehension. The benefits of this approach are multi-faceted:  

\begin{itemize}
    \item \textbf{Coverage of Key Topics:} The game includes questions on significant figures, unit conversions, basics of matter, story problems, and fun chemistry facts. This ensures that students revisit each major concept from the chapters, from understanding the proper number of digits in measurements to applying unit conversions in practical scenarios, such as calculating the volume of a liquid medication or determining the mass of a sample in a lab experiment.
    
    \item \textbf{Formative Assessment:} By participating in the game, students gain immediate feedback on their understanding. For example, if a team struggles to answer a question about converting grams to moles, they recognize that additional practice is needed before moving on to more advanced topics. This self-assessment mirrors real-life learning, where recognizing gaps in knowledge is the first step toward improvement.
    
    \item \textbf{Collaborative Learning:} Students work in teams, encouraging discussion, debate, and collective problem-solving. For instance, when faced with a story problem about mixing solutions, team members must communicate and decide which approach to take. This mirrors real-life scientific work, where collaboration and peer feedback often lead to better outcomes than working in isolation.
    
    \item \textbf{Preparation and Study Habits:} The game highlights the importance of preparation. Students who review class notes, practice calculations, and understand core concepts perform better. This reflects real-life situations, such as preparing for lab experiments or exams, where consistent preparation leads to confidence and success.
    
    \item \textbf{Engagement and Motivation:} Incorporating game elements makes learning enjoyable, helping students stay motivated. For example, earning points for correctly answering a question about chemical reactions can make abstract concepts more tangible, similar to using simulations or interactive apps to visualize chemical processes in the real world.
    
    \item \textbf{Reflection on Learning:} After the game, students can reflect on their performance. They can ask themselves questions like: Which topics were easy, which were challenging, and what strategies helped my team succeed? This reflection encourages metacognition, the real-world skill of thinking about one’s own learning process, which is crucial in both academic and professional settings.
\end{itemize} \\
\hline
\end{tabular}
\newpage

\noindent
\section*{\underline{Activity 2}}
\begin{tabular}{|p{4.5cm}|p{9.5cm}|}
\hline
\textbf{Collaborative Learning Techniques} & Fishbowl (many), Jigsaw (few)\\
\hline
\textbf{Learning Strategy} &  Boardwork Model (modified)\\
\hline
\textbf{Time} & 12:30---12:50\\
\hline
\textbf{Materials} & Colored whiteboard markers, a whiteboard, partitioned questions (to display multiple to multiple groups at a time).\\
\hline
\textbf{Lecture/Textbook/Old Plans/Slide References} & Chapter 1: Basics of Matter, Chapter 2: Measurements\\
\hline
\end{tabular}

\vspace{0.5cm}
\noindent
\begin{tabular}{|p{0.9\linewidth}|}
\hline
{\large\textbf{Definitions}}
\begin{enumerate}
   \item\textbf{Fishbowl:} The group is separated into an "inside" and an "outside" the fishbowl. Students inside the fishbowl will participate in the activity or discussion. Students outside the fishbowl will listen to the discussion and observe the students engage in the activity. After the fishbowl, students inside are given the opportunity to reflect on what they observed, provide feedback to those isnide, and can contribute more to the discussion.
   \item\textbf{Jigsaw:} Jigsaws, when used properly, make the group as a whole dependent upon all of them in subgroups. Each group provides a peice of the puzzle. Group members are broken into smaller groups. Each small group works on some aspect of the same problem, question, or issue. They then share their part of the puzzle with the large group.
   \item\textbf{Boardwork Model:} This is a method of organizing boardwork in order to fascilitate an undertsanding of problem-solving strategies. The baord should be divdied into 4 sections: 1.) Prerequisite Knowledge, 2.) Mathematical Steps, 3.) Narrative of the Steps, 4.) Additional Sample Problem. Encourage one student to fill out section 1 on the board. Then encourage 2 students to simultaneously complete section 2 and 3 on the board. Lastly, have another student complete section 4. Encourage the students to use this model when studying outside of the SI session.
\end{enumerate}
\\ \hline
\end{tabular}
\newpage
{\large\textbf{Instructions (please provide a\hl{DETAILED} activity breakdown below (and/or attached)}}\\
\begin{tabular}{|p{0.9\linewidth}|}
\hline
\begin{enumerate}
\item\textbf{Objective:}
Guide students to work in groups of at least 2 to solve chemistry-based problems collaboratively on a whiteboard.
\end{enumerate}

\hline  
\textbf{Setup}  
\begin{itemize}  
   \item Arrange the classroom so multiple groups can work at the whiteboard simultaneously.  
   \item Provide each group with markers and a section of the whiteboard.  
   \item Prepare a set of chemistry problems in advance.  
   \item Optionally, provide reference tables or hints for problem-solving to keep groups on track.  
\end{itemize}  

\textbf{Running the Activity}  
\begin{enumerate}  
   \item Introduce the activity and explain that groups of at least 2 will calculate problem results together.  
   \item Present the first chemistry problem to all groups.  
   \item Groups collaborate to solve the problem on the whiteboard.  
   \item Rotate through problems, allowing each group to contribute their calculations and drawings.  
   \item Monitor group progress, providing hints or clarifications as needed without giving away solutions.  
   \item Encourage discussion and peer teaching within groups to reinforce chemistry concepts.  
   \item Continue until all problems are solved and the final picture is completed.  
\end{enumerate}  

\textbf{Facilitator Role}  
\begin{itemize}  
   \item Ensure groups are on task and collaborating effectively.  
   \item Clarify instructions and answer conceptual questions without solving problems for the students.  
   \item Track timing to ensure the activity progresses smoothly.  
   \item Circulate to observe problem-solving strategies, noting areas where students struggle for future review.  
   \item Celebrate the completion of the final picture and debrief with the class on the chemistry concepts applied.  
\end{itemize}  
\\ \hline  
\end{tabular}  

\newpage
\noindent
{\large\textbf{Checking for Understanding (questions & answers)}}\\[0.2em]
\begin{tabular}{|p{0.9\linewidth}|}
\hline
\textbf{Questions:}
\begin{enumerate}
\item Why does the boardwork model provide a calmer, low-pressure environment for students?  
\item How does the centralized canvas in the boardwork model help reduce student anxiety about mistakes?  
\item In what ways can students approach chemistry problems when using the boardwork model?  
\item What types of chemistry tasks can students work through step by step with this model?  
\item How do peers act as resources in the boardwork model?  
\item How does peer explanation and feedback reinforce understanding of chemistry concepts?  
\item Within a Supplemental Instruction (SI) context, how does the boardwork model encourage participation?  
\item What benefits do students gain from multiple perspectives on the same chemistry problem?  
\item How does the shared board help students practice applying theoretical knowledge in a visual way?  
\item What role does the SI leader play when using the boardwork model during a session?  
\item How does this model help the SI leader identify common misconceptions?  
\item In what way does the boardwork model make SI sessions more effective for learners of varied skill levels?  
\end{enumerate} \\
\hline  
\end{tabular}

\newpage
\noindent
\section*{\underline{Closer}}
\begin{tabular}{|p{4.2cm}|p{9.8cm}|}
\hline
\textbf{Collaborative Learning Techniques} & Group Survey\\
\hline
\textbf{Learning Strategy} & Brain dump\\
\hline
\textbf{Time} & 12:50---12:55\\
\hline
\textbf{Materials} & Index cards, writing utensils\\
\hline
\textbf{Lecture/Textbook/Old Plans/Slide References} & Chapter 1: Basics of Matter, Chapter 2: Measurements\\
\hline
\end{tabular}

\vspace{0.5cm}
\noindent
\begin{tabular}{|p{0.9\linewidth}|}
\hline
{\large\textbf{Definitions}}
\begin{enumerate}
   \item\textbf{Group survey:} Each group member is surveyed to discover their position on an issue, problem or topic. This process ensures that each member of the group is allowed to offer or state their point of view.
   \item\textbf{Brain dump:} This is a great closer for students to sum up their knowledge on the topic(s) from the session. Have the students take out a blank piece of paper /Google Doc and write/type("dump") everything they know relevant from the discussion that day. The SI Leader can look at their work to assess the students' understanding of the course material and have them compare it to each other. It is important to at least discuss with the students so the SI Leader knows the students' comfort levels with the material. 
\end{enumerate}
\\ \hline
\end{tabular}

\newpage
\noindent
{\large\textbf{Instructions (please provide a\hl{DETAILED} activity breakdown below (and/or attached}}

\begin{tabular}{|p{0.9\linewidth}|}
\hline
\begin{enumerate}[itemsep=0.3em]
   \item Students are put into small groups and are asked to write down their confidence on a single notecard without names.
   \item The cards are evidence of the group's confidence in the content.
\end{enumerate}
\\ \hline
\end{tabular}

\vspace{0.5cm}
\noindent
{\large\textbf{Content (fully worked out problems \& answers)}

\begin{tabular}{|p{0.9\linewidth}|}
\hline
\begin{itemize}
   \item No new worked problems; brain dump.
\end{itemize}
\\ \hline
\end{tabular}

\vspace{0.5cm}
\noindent
{\large\textbf{Checking for Understanding (questions \& answers) continued...}

\begin{tabular}{|p{0.9\linewidth}|}
\hline
\begin{itemize}
   \item Review confidence levels
   \item Interpret the variety of confidence and their distribution in groups.
\end{itemize}
\\ \hline
\end{tabular}
\newpage
\section*{Plan Checklist}

\begin{tabular}{|p{0.9\linewidth}|}
\hline
\begin{itemize}
\item[$\Box$] Variety of Collaborative Learning Techniques and Learning Strategies
\item[$\Box$] Chapter/slide\# where information is coming from
\item[$\Box$] Included timestamps (and buffer time as needed)
\item[$\Box$] Details on how leader will break up students
\item[$\Box$] Checking for Understanding questions
\item[$\Box$] Fully worked out problems and examples for each activity
\item[$\Box$] Necessary links required for online implementation (if applicable)
\item[$\Box$] Source material correctly cited (if applicable)
\item[$\Box$] Plan is uploaded to Teams
\end{itemize}
\\ \hline
\end{tabular}
\newpage

\end{document}
